
\documentclass[11pt]{article}

% --- page + fonts ---
\usepackage[margin=1in]{geometry}
\usepackage{microtype}
\usepackage[T1]{fontenc}

% --- math ---
\usepackage{amsmath,amssymb,amsthm,mathtools}
\usepackage{physics}

% --- misc ---
\usepackage{cite}
\usepackage{hyperref}
\usepackage{enumitem}
\usepackage{booktabs}
\usepackage{xcolor}
\hypersetup{pdfborder={0 0 0}}
% --- camera-ready linebreak tolerance ---
\setlength{\emergencystretch}{3em}


% --- theorem environments ---
\newtheorem{definition}{Definition}[section]
\newtheorem{theorem}[definition]{Theorem}
\newtheorem{lemma}[definition]{Lemma}
\newtheorem{proposition}[definition]{Proposition}
\newtheorem{corollary}[definition]{Corollary}
\theoremstyle{remark}
\newtheorem{remark}[definition]{Remark}

% --- macros (notation deliberately minimal) ---
\newcommand{\SM}{\ensuremath{\mathrm{SM}}}
\newcommand{\EFT}{\ensuremath{\mathrm{EFT}}}
\newcommand{\CC}{\ensuremath{\mathrm{CC}}}
\newcommand{\NC}{\ensuremath{\mathrm{NC}}}
\newcommand{\EM}{\ensuremath{\mathrm{EM}}}
\newcommand{\SU}{\ensuremath{\mathrm{SU}}}
\newcommand{\U}{\ensuremath{\mathrm{U}}}
\newcommand{\cO}{\mathcal{O}}
\newcommand{\Skin}{\mathcal{S}}
\newcommand{\Wit}{\mathcal{W}}
\newcommand{\Load}{\mathcal{L}}
\newcommand{\Ad}{\ensuremath{\mathsf{Ad}}}
\newcommand{\Jarlskog}{\mathcal{J}}
\DeclareMathOperator{\diag}{diag}

\title{\textbf{Structura Ex Necessitate Standardis Modelis}\\
\large A standalone forced-derivation ledger for Standard-Model bookkeeping\\
\large under admissible redescription and transport closure}
\author{Amos Jay Maley}
\date{\today}

\begin{document}
\maketitle

\begin{abstract}
This paper is \emph{standalone} and \emph{audit-grade}. It does not propose new dynamics or ontology;
it audits the internal closure of a fixed renormalizable chiral gauge load
(Definition~\ref{def:declared-load}) under a single governance requirement:
\emph{standing} is what survives admissible redescription together with explicit, scope-indexed regime
transport. Every collapse is exhibited as a \emph{transport certificate}
(Definition~\ref{def:certificate}): the skins are named, the transport map is written, the invariant
witness set is computed, and witness rank is compared to apparent parameter rank.
Low-energy sector results include charged-current matching to $G_F$, electromagnetic matching to $e$,
Yukawa--mass equivalences, $\bar\theta$ as the unique CP--odd QCD invariant, oscillation invariants, and
the longitudinal--Goldstone amplitude-class collapse. Discrete-skeleton forcing is then made explicit
under the minimal chiral skeleton with one Higgs doublet: anomaly cancellation forces the hypercharge
class and guarantees existence of the charged Yukawa operators; hypercharge normalization is shown
redundant; tree-level FCNC are excluded. We then define an admissibility envelope and prove a
\emph{relative} exhaustion theorem: within the envelope and declared load, no further equivalence
completion exists without introducing new load, leaving the scope, or invoking a standing-failure mode.
Finally, we state the witness normal form for the full renormalizable scope and certify which
parameters remain irreducible inputs and why their numerical values are not internally derivable
without breaking redescription stability. We also prove a meta-closure result: under minimal coherence, the standing quotient and the admissibility groupoid are unique up to isomorphism and admit no proper same-scope standing-preserving extension.
\end{abstract}

\tableofcontents

\section{Scope and non-negotiables}

\subsection{Method posture}
This paper is eliminative. It introduces:
\begin{itemize}[leftmargin=2.2em]
\item no new primitives, no new ontology, and no new anchoring claims;
\item no compensating fields, auxiliary sectors, additional $U(1)$ factors, or exotic charge lattices;
\item no probabilistic/typicality repairs, no semantic standing import, and no privileged representation.
\end{itemize}
All results are obtained by explicit calculation plus one governance rule: \emph{standing is what
survives admissible redescription and admissible transport}. Anything else is redundancy or standing failure.

\subsection{Declared load}
We fix the following \emph{declared load} for the entire paper.

\begin{definition}[Declared load]\label{def:declared-load}
The declared load $\Load$ consists of:
\begin{enumerate}[label=(\roman*),leftmargin=2.2em]
\item the renormalizable $\SU(3)_c\times \SU(2)_L\times \U(1)_Y$ gauge framework;
\item minimal chiral fermion content in the usual representation skeleton (three generations allowed but
not required for the algebraic derivations);
\item a single electroweak Higgs doublet $H$ (with conjugate $\tilde H \equiv i\sigma^2 H^\ast$ available
as usual);
\item no additional Higgs multiplets, no vectorlike patch pairs, no extra abelian factors, and no
\emph{ad hoc} patch operators (non-gauge-invariant or higher-dimensional) introduced as independent
primitives; higher-dimensional operators may appear only as \emph{transport images} (EFT skins) of the
renormalizable load.
\end{enumerate}
\end{definition}

\begin{remark}[What is and is not derived]\label{rem:not-derive-gauge}
Definition~\ref{def:declared-load} fixes the renormalizable chiral gauge skeleton whose internal closure is
audited. This paper does \emph{not} attempt to derive the gauge group or the representation skeleton from
scratch; doing so would be a different project with a different scope.
The present derivations concern (i) collapse of \emph{apparent} parameter freedoms under scope-indexed
transport, and (ii) closure of \emph{same-scope} modifications relative to the fixed load.
\end{remark}

\begin{remark}[Fail-closed discipline]\label{rem:fail-closed}
\emph{Inadmissible} is used in the fail-closed sense: if a purported move cannot be made
standing-preserving under admissible redescription and explicit transport inside $\Load$, then it is
rejected as either (a) redundancy (no standing), or (b) standing failure (requiring hidden load,
representational privilege, illicit scope transport, or compensating structure).
\end{remark}



\subsection{Declared load is an envelope, not a hidden derivation anchor}\label{sec:load-envelope}
A frequent misunderstanding in ``uniqueness'' arguments is to treat a fixed load choice as a covert
assumption smuggled in to force the desired answer. Here we make the opposite move: we \emph{type} the
claim. The paper does \emph{not} assert that no other loads are logically possible. It asserts that,
\emph{relative to the scope and declared load}, there is a unique \emph{minimal standing-preserving
interior}, and that purported ``extensions'' that claim to remain in the same scope are either
redundant or scope-changing.

\begin{definition}[Scope-preserving vs.\ scope-changing modification]\label{def:scope-preserving}
Fix a scope $S$ (Definition~\ref{def:scope}).
A proposed modification $\Delta$ of the bookkeeping is:
\begin{enumerate}[label=(\roman*),leftmargin=2.2em]
\item \emph{scope-preserving} if it claims to reproduce exactly the same invariant witness set for $S$
(up to admissible redescription) without introducing additional standing-relevant witnesses;
\item \emph{scope-changing} otherwise (it enlarges/changes the witness set, the admissibility envelope,
or the regime content).
\end{enumerate}
\end{definition}

\begin{proposition}[Added-load dichotomy]\label{prop:added-load}
Let $S$ be a fixed scope and $\Load$ the declared load.
Consider adding extra fields/operators/sectors $\Delta\Load$.
If $\Delta\Load$ introduces \emph{no} new invariant witnesses for $S$, then $\Delta\Load$ is redundancy
for $S$ (it is removable by admissible redescription/transport, e.g.\ integrating out).
If $\Delta\Load$ \emph{does} introduce new invariant witnesses for $S$, then it is a scope change and
therefore not a nontrivial \emph{scope-preserving} extension of the \SM interior.
\end{proposition}

\begin{proof}
If $\Delta\Load$ introduces no new invariant witnesses for $S$, then by Definition~\ref{def:tensor-load}
it contributes no standing-relevant data for $S$ and is removable without changing witness content.
If $\Delta\Load$ introduces new invariant witnesses, then by Definition~\ref{def:scope-preserving} the
scope has changed: the modified bookkeeping describes strictly more (or different) standing content.
In neither case does $\Delta\Load$ produce a genuine same-scope enlargement of standing content.
\end{proof}


\begin{remark}[What is \emph{not} being claimed]\label{rem:not-bsm-ban}
Proposition~\ref{prop:added-load} is not a ``BSM ban.'' It is a typing statement:
\emph{inside} the present scope and declared load, additional structure is either redundant or a scope
change. Scope changes may be physically motivated; they are simply different projects and are not
adjudicated here.
\end{remark}
\section{Governance layer: admissibility, transport, witnesses}

\subsection{Minimal coherence forces the admissibility interface}\label{sec:min-coherence}

The governance primitives in this section are not decorative. They are\footnote{In the same sense that
Myhill--Nerode equivalence is not ``chosen'' but forced once one demands a canonical presentation of
future-distinguishability classes.} forced by minimal coherence requirements for any scope-indexed
bookkeeping practice.

\begin{definition}[Minimal coherence requirements]\label{def:coherence}
Fix a scope $S$.
A bookkeeping practice for $S$ is \emph{minimally coherent} if it satisfies:
\begin{enumerate}[label=(\textbf{MC\arabic*}),leftmargin=2.8em]
\item \textbf{Stable reference.} Empirical referents tracked in $S$ must remain identifiable across lawful
reformulations; the identity of ``the same'' referent may not depend on a privileged internal
parameterization.
\item \textbf{Lawful redescription freedom.} Multiple internal representations may be used (field variables,
bases, gauge-fixings, parameterizations), and such redescriptions are permitted so long as they preserve
$S$-predictions.
\item \textbf{Irreversible commitment.} Once an empirical record is fixed (a completed inference, stored data),
it cannot be retroactively reindexed by choosing a different internal representation; coherence must be
maintained under irreversible use of records.
\item \textbf{Non-triviality.} The bookkeeping makes at least one prediction that can distinguish some admissible
cases inside $S$ (otherwise all ``states'' collapse vacuously).
\end{enumerate}
\end{definition}

\begin{remark}[MC1 is a standing postulate, not an ontic prohibition]\label{rem:mc1-standing}
The minimal-coherence axioms are governance constraints for \emph{standing} in a fixed scope $S$.
In particular, \textbf{MC1} says that within the same-scope interior we do not treat distinctions with no admissible
$S$-witness as identity-bearing.
This is \emph{not} the metaphysical claim that nature contains no additional structure; it is the bookkeeping
discipline that such additional structure, if asserted, must appear as added load or enlarged scope and hence be
typed as a scope change (Proposition~\ref{prop:added-load}).
\end{remark}


\begin{theorem}[Minimal coherence forces the standing quotient and admissibility groupoid]\label{thm:coherence-groupoid}
Assume minimal coherence (Definition~\ref{def:coherence}) for a fixed scope $S$.
Then:
\begin{enumerate}[label=(\roman*),leftmargin=2.2em]
\item The only identity-bearing objects for $S$ are standing equivalence classes $[\mathbf{p}]\in P(\Skin)/\!\sim_S$
(Definitions~\ref{def:standing-equiv}--\ref{def:standing-quotient}). In particular, any attempt to treat two
parameter points $\mathbf{p}\sim_S\mathbf{p}'$ as distinct standing states introduces representational
privilege or hidden unwitnessed structure (Definition~\ref{def:failure}).
\item Any claim that two skins $\Skin_A,\Skin_B$ represent the \emph{same} scope content requires transports that
are transport-complete for $S$ (Definition~\ref{def:transport-complete}); equivalently, the induced maps on
standing quotients must be injective, and for mutual equivalence bijective.
\item The class of skins equipped with transport-complete admissible transports that admit inverses on standing
quotients forms a \emph{groupoid} $\Ad_S$ (objects: skins for $S$; morphisms: transports inducing bijections on
standing quotients). Standing is functorial: every morphism in $\Ad_S$ induces a bijection of standing state
spaces.
\end{enumerate}
\end{theorem}

\begin{proof}
\textbf{(i)} Let $\Skin$ be any skin used for scope $S$. By (MC2), lawful redescriptions are permitted provided they
preserve $S$-predictions. Therefore any internal parameter value that does not change $S$-predictions cannot be
privileged as identity-grounding without violating (MC1): doing so would make referent identity depend on an
arbitrary representational choice.

Formally, define $\mathbf{p}\sim_S\mathbf{p}'$ exactly as in Definition~\ref{def:standing-equiv}: equality of all
$S$-predictions within the declared error discipline. If $\mathbf{p}\sim_S\mathbf{p}'$ but one insists they are
distinct standing states, then the distinction has \emph{no} $S$-observable witness.
Maintaining the distinction therefore requires either a privileged parameterization (violating MC1/MC2), or
extra unwitnessed structure imported to ``remember'' which point one is in (hidden load), which is precisely a
standing-failure mode (Definition~\ref{def:failure}).
Thus coherence forces identification of standing states with equivalence classes $[\mathbf{p}]$.

\textbf{(ii)} Consider two skins $\Skin_A,\Skin_B$ both purported to describe scope $S$.
If there is no transport-complete map, then there exist $\mathbf{p},\mathbf{p}'$ in one skin that are
$S$-distinguishable but whose images in the other skin cannot be distinguished by any invariant witness there.
Then the identity of the $S$-referent depends on which skin is used, contradicting (MC1) and (MC3).
Therefore any coherence claim of ``same scope content'' requires transport completeness.
Injectivity/bijectivity on standing quotients is exactly this requirement in quotient form.

\textbf{(iii)} Define $\Ad_S$ to have as objects those skins used for $S$ and as morphisms the admissible transports
that induce bijections on standing quotients (equivalently, are transport-complete and admit inverses up to
$\sim_S$). Composition of such morphisms is again transport-complete and induces bijections on quotients; identity
morphisms are given by trivial transports. Existence of inverses on quotients makes this a groupoid.
Functoriality of standing is then immediate: each morphism induces a bijection on $P(\Skin)/\!\sim_S$ by
construction; Lemma~\ref{lem:quotient-iso} provides the explicit statement.
\end{proof}



\subsection{Meta-closure: uniqueness and maximality}\label{sec:meta-closure}
\paragraph{Referee Proof Obligations.}
To refute the meta-closure results of Section~\ref{sec:meta-closure}, a critic must do \emph{at least one} of the following:
(i) explicitly deny at least one Minimal Coherence axiom (MC1--MC4), stating which condition of stable reference, irreversibility, lawful redescription, or non-trivial inference is rejected;
(ii) exhibit a concrete standing presentation $(X,r)$ satisfying (MC1--MC4) that reproduces all $S$-predictions yet is not isomorphic (via the uniquely induced bijection) to the standing quotient $P(\mathrm{Skin})/\!\sim_S$;
or (iii) produce a same-scope morphism that introduces a new invariant witness without violating transport completeness, minimal coherence, or declared load.
Absent one of (i)--(iii), objections reduce to scope change, load augmentation, or representational privilege, each already classified.

The previous subsection established that minimal coherence forces the quotient notion of standing and the
groupoid form of same-scope equivalence. A hostile reader may still attempt a meta-level escape:
``perhaps there exists some \emph{other} equally coherent identity/transport scheme.''
Here we close that route in the same way as a Myhill--Nerode argument closes ``alternative'' state
presentations: the standing quotient is not one option among many; it is the unique standing state space
compatible with scope-indexed prediction invariance and stable reference.

\begin{definition}[Standing presentation for a scope]\label{def:standing-presentation}
Fix a scope $S$ and a skin $\Skin$ with parameter space $P(\Skin)$.
A \emph{standing presentation} for $S$ is a triple $(X,r,\mathrm{Pred}_X)$ where:
\begin{enumerate}[label=(\roman*),leftmargin=2.2em]
\item $X$ is a set intended to serve as the standing state space for $S$;
\item $r:P(\Skin)\to X$ is a surjection (``representation map'');
\item for every $\mathcal{O}\in S$ there is a prediction map
$\mathrm{Pred}_X(\mathcal{O};\cdot):X\to\mathbb{R}$ such that
\begin{equation}\label{eq:standing-presentation}
\mathrm{Pred}_{\Skin}(\mathcal{O};\mathbf{p})=\mathrm{Pred}_X(\mathcal{O};r(\mathbf{p}))
\qquad (\forall\,\mathbf{p}\in P(\Skin),\ \forall\,\mathcal{O}\in S).
\end{equation}
\end{enumerate}
We call the presentation \emph{sound} if $\mathbf{p}\sim_S\mathbf{p}'\Rightarrow r(\mathbf{p})=r(\mathbf{p}')$,
and \emph{complete} if $r(\mathbf{p})=r(\mathbf{p}')\Rightarrow \mathbf{p}\sim_S\mathbf{p}'$
(Definition~\ref{def:standing-equiv}).
\end{definition}

\begin{lemma}[Completeness of standing presentations]\label{lem:standing-complete}
Let $(X,r,\mathrm{Pred}_X)$ be a standing presentation for scope $S$ in the sense of
Definition~\ref{def:standing-presentation}.
Then the presentation is \emph{complete}:
\[
r(\mathbf{p})=r(\mathbf{p}')\ \Longrightarrow\ \mathbf{p}\sim_S \mathbf{p}'.
\]
\end{lemma}

\begin{proof}
Assume $r(\mathbf{p})=r(\mathbf{p}')$.
Then for every $\mathcal{O}\in S$, the factorization condition \eqref{eq:standing-presentation} gives
\[
\mathrm{Pred}_{\Skin}(\mathcal{O};\mathbf{p})
=\mathrm{Pred}_X(\mathcal{O};r(\mathbf{p}))
=\mathrm{Pred}_X(\mathcal{O};r(\mathbf{p}'))
=\mathrm{Pred}_{\Skin}(\mathcal{O};\mathbf{p}').
\]
Hence $\mathbf{p}\sim_S\mathbf{p}'$ by Definition~\ref{def:standing-equiv}.
\end{proof}

\begin{lemma}[Soundness forced by stable reference]\label{lem:standing-sound}
Assume minimal coherence for scope $S$ (Definition~\ref{def:coherence}).
Let $(X,r,\mathrm{Pred}_X)$ be a standing presentation for $S$ in the sense of
Definition~\ref{def:standing-presentation}.
Then the presentation is \emph{sound}:
\[
\mathbf{p}\sim_S\mathbf{p}'\ \Longrightarrow\ r(\mathbf{p})=r(\mathbf{p}').
\]
\end{lemma}

\begin{proof}
Assume $\mathbf{p}\sim_S\mathbf{p}'$.
Then for every $\mathcal{O}\in S$,
$\mathrm{Pred}_{\Skin}(\mathcal{O};\mathbf{p})=\mathrm{Pred}_{\Skin}(\mathcal{O};\mathbf{p}')$.
By \eqref{eq:standing-presentation} this implies
$\mathrm{Pred}_X(\mathcal{O};r(\mathbf{p}))=\mathrm{Pred}_X(\mathcal{O};r(\mathbf{p}'))$ for all $\mathcal{O}\in S$.

If $r(\mathbf{p})\neq r(\mathbf{p}')$, then $X$ distinguishes two putative ``standing states'' that are
indistinguishable by \emph{every} admissible $S$-observable. Maintaining such an unwitnessed distinction as
identity-bearing violates stable reference (MC1) under the fail-closed discipline: it either (i) privileges the
particular $r$-labeling as identity-grounding (representational privilege), or (ii) imports extra unwitnessed
structure to remember which label one is in (hidden load). Both are standing-failure modes
(Definition~\ref{def:failure}). Therefore $r(\mathbf{p})=r(\mathbf{p}')$ and the presentation is sound.
\end{proof}

\begin{lemma}[Quotient universal property]\label{lem:quotient-universal}
Let $q:P(\Skin)\to P(\Skin)/\!\sim_S$ be the quotient map.
For any set $Y$ and any function $F:P(\Skin)\to Y$ that is constant on $\sim_S$-classes
(i.e.\ $\mathbf{p}\sim_S\mathbf{p}'\Rightarrow F(\mathbf{p})=F(\mathbf{p}')$),
there exists a unique function $\widetilde F:P(\Skin)/\!\sim_S\to Y$ such that
\[
F=\widetilde F\circ q.
\]
\end{lemma}

\begin{proof}
Define $\widetilde F([\mathbf{p}]):=F(\mathbf{p})$.
This is well-defined precisely because $F$ is constant on $\sim_S$-classes.
Uniqueness follows because $q$ is surjective: if $F=\widehat F\circ q$ also, then for any $[\mathbf{p}]$,
$\widehat F([\mathbf{p}])=\widehat F(q(\mathbf{p}))=F(\mathbf{p})=\widetilde F(q(\mathbf{p}))=\widetilde F([\mathbf{p}])$.
\end{proof}

\begin{theorem}[Uniqueness of the standing quotient]\label{thm:standing-quotient-unique}
Assume minimal coherence for scope $S$ (Definition~\ref{def:coherence}).
Let $(X,r,\mathrm{Pred}_X)$ be any standing presentation for $S$ in the sense of
Definition~\ref{def:standing-presentation}.
Then $r$ is automatically sound and complete and induces a unique bijection
\[
\phi:\ P(\Skin)/\!\sim_S\ \longrightarrow\ X,\qquad \phi([\mathbf{p}])=r(\mathbf{p}).
\]
In particular, the standing quotient $P(\Skin)/\!\sim_S$ is the \emph{unique} standing state space for $S$
(up to unique isomorphism) compatible with stable reference and scope-indexed prediction invariance.
Equivalently: there is no non-isomorphic coherent identity framework for $S$ that reproduces all $S$-predictions.
\end{theorem}

\begin{proof}
By Lemmas~\ref{lem:standing-complete} and \ref{lem:standing-sound}, the presentation is complete and sound.
Define $\phi([\mathbf{p}]):=r(\mathbf{p})$.

Soundness implies $\phi$ is well-defined (it is constant on $\sim_S$-classes).
Completeness implies $\phi$ is injective: if $\phi([\mathbf{p}])=\phi([\mathbf{p}'])$ then
$r(\mathbf{p})=r(\mathbf{p}')$ and Lemma~\ref{lem:standing-complete} gives $\mathbf{p}\sim_S\mathbf{p}'$, hence
$[\mathbf{p}]=[\mathbf{p}']$.
Surjectivity holds because $r$ is surjective.

Uniqueness of $\phi$ is exactly the quotient universal property (Lemma~\ref{lem:quotient-universal}):
$r$ is constant on $\sim_S$-classes (soundness), so it factors through $q$ uniquely as $r=\phi\circ q$.
\end{proof}

\paragraph{Realist coherence variant: continuation-separability as normal form.}
A recurring hostile move is to claim that \textbf{MC1} is ``operationalist'': a realist may allow additional
structure not distinguished by the base observable list. The meta-closure results above do not require that
stronger metaphysical denial. What is required for \emph{identity-bearing standing} is that alleged identity
distinctions be \emph{separable by some admissible continuation} in the declared scope.
We therefore make the Myhill--Nerode analogy literal by defining admissible histories $\mathcal{H}_S$ as the
closure of admitted tests under admissible interventions, and by exhibiting a canonical \emph{reduction} of any
standing presentation whose standing states are continuation-separated. In this sense \textbf{CR-Sep} is not an
additional metaphysical axiom; it is the minimality (no redundant labels) normal form of any same-scope standing
presentation.

\begin{definition}[Admissible histories $\mathcal{H}_S$ as closure under composition]\label{def:histories}
Fix a scope declaration $S=(S_0,\mathcal{U}_S)$ and a skin $\Skin$ with parameter space $P(\Skin)$, where $S_0$ is the admitted observable family and $\mathcal{U}_S(\Skin)$ is the monoid of admissible interventions/updates.
Let $\mathcal{U}_S(\Skin)\subseteq \mathrm{End}(P(\Skin))$ denote the monoid of admissible
interventions/updates in scope $S$ (it contains $\mathrm{id}$ and is closed under composition).
For each admitted observable $\mathcal{O}\in S_0$ define the associated \emph{test map}
$t_{\mathcal{O}}:P(\Skin)\to\mathbb{R}$ by
\[
t_{\mathcal{O}}(\mathbf{p}) := \mathrm{Pred}_{\Skin}(\mathcal{O};\mathbf{p}).
\]
Define the set of admissible \emph{histories/continuations} to be the closure of the primitive tests under
precomposition by admissible interventions:
\[
\mathcal{H}_S(\Skin)\ :=\ \{\,t_{\mathcal{O}}\circ u\ :\ \mathcal{O}\in S_0,\ u\in\mathcal{U}_S(\Skin)\,\}.
\]
Equivalently we may index a history by a pair $h=(u,\mathcal{O})$ with $u\in\mathcal{U}_S(\Skin)$ and
$\mathcal{O}\in S_0$ and write
\[
\mathrm{Pred}_{\Skin}(h;\mathbf{p})\ :=\ \mathrm{Pred}_{\Skin}(\mathcal{O};u(\mathbf{p})).
\]
Define the \emph{continuation signature} map
\[
\Sigma_{\Skin}:\ P(\Skin)\longrightarrow \mathbb{R}^{\mathcal{H}_S(\Skin)},\qquad
\Sigma_{\Skin}(\mathbf{p}) := \big(\mathrm{Pred}_{\Skin}(h;\mathbf{p})\big)_{h\in\mathcal{H}_S(\Skin)}.
\]
When $\Skin$ is fixed we write $\mathcal{H}_S:=\mathcal{H}_S(\Skin)$.
We call a scope $S$ \emph{history-saturated} (relative to $\mathcal{U}_S$) if it already contains all composite
tests $t_{\mathcal{O}}\circ u$ as admitted observables.
In that case Definition~\ref{def:standing-equiv} may be read as quantifying over $\mathcal{H}_S$, and the present
notation makes the Myhill--Nerode ``future-distinguishability'' content explicit rather than rhetorical.
\end{definition}

\begin{lemma}[Kernel equivalence is the maximal update-stable refinement of primitive prediction equivalence]
\label{lem:kernel-equivalence}
Fix a scope declaration $S=(S_0,\mathcal{U}_S)$ and a skin $\Skin$.
Define the \emph{primitive} (snapshot) prediction-equivalence relation $\sim_{0}$ on $P(\Skin)$ by
\[
\mathbf{p}\sim_{0}\mathbf{p}' \quad\Longleftrightarrow\quad
\mathrm{Pred}_{\Skin}(\mathcal{O};\mathbf{p})=\mathrm{Pred}_{\Skin}(\mathcal{O};\mathbf{p}')
\ \ \forall\,\mathcal{O}\in S_0.
\]
Define the \emph{continuation} (kernel) equivalence relation $\sim_{\ker}$ by
\[
\mathbf{p}\sim_{\ker}\mathbf{p}' \quad\Longleftrightarrow\quad
\Sigma_{\Skin}(\mathbf{p})=\Sigma_{\Skin}(\mathbf{p}'),
\]
where $\Sigma_{\Skin}$ is the continuation-signature map of Definition~\ref{def:histories}.
Equivalently, $\mathbf{p}\sim_{\ker}\mathbf{p}'$ iff
$\mathrm{Pred}_{\Skin}(\mathcal{O};u(\mathbf{p}))=\mathrm{Pred}_{\Skin}(\mathcal{O};u(\mathbf{p}'))$
for all $\mathcal{O}\in S_0$ and all $u\in\mathcal{U}_S(\Skin)$.

Then:
\begin{enumerate}[label=(\roman*),leftmargin=2.2em]
\item $\sim_{\ker}$ is a right-congruence for admissible update composition:
if $\mathbf{p}\sim_{\ker}\mathbf{p}'$ then $u(\mathbf{p})\sim_{\ker}u(\mathbf{p}')$ for all $u\in\mathcal{U}_S(\Skin)$.
\item $\sim_{\ker}$ is the \emph{largest} right-congruence contained in $\sim_{0}$.
That is: if $\approx$ is any equivalence relation on $P(\Skin)$ such that
(a) $\mathbf{p}\approx \mathbf{p}'\Rightarrow \mathbf{p}\sim_{0}\mathbf{p}'$ and
(b) $\mathbf{p}\approx \mathbf{p}'\Rightarrow u(\mathbf{p})\approx u(\mathbf{p}')$ for all $u\in\mathcal{U}_S(\Skin)$,
then $\mathbf{p}\approx \mathbf{p}'\Rightarrow \mathbf{p}\sim_{\ker}\mathbf{p}'$.
\item Consequently, if there exist $\mathbf{p},\mathbf{p}'$ with $\mathbf{p}\sim_{0}\mathbf{p}'$ but $\mathbf{p}\not\sim_{\ker}\mathbf{p}'$
(kernel-inequivalent but snapshot-equivalent), then no quotient that identifies $\mathbf{p}$ and $\mathbf{p}'$ can admit
a well-defined induced action of $\mathcal{U}_S(\Skin)$ by $[\mathbf{p}]\mapsto [u(\mathbf{p})]$.
Under (MC3) irreversible commitment (Definition~\ref{def:coherence}), this ambiguity cannot be repaired post hoc by
retroactive representative selection; any coherent identity notion compatible with composed updates must refine to $\sim_{\ker}$.
In a history-saturated scope, $\sim_{\ker}$ coincides with $\sim_S$ (Definition~\ref{def:histories}), so identity collapses to the standing quotient.
\end{enumerate}
\end{lemma}

\begin{proof}
\textbf{(i)} Assume $\mathbf{p}\sim_{\ker}\mathbf{p}'$, i.e.\ $\Sigma_{\Skin}(\mathbf{p})=\Sigma_{\Skin}(\mathbf{p}')$.
Fix any $u_0\in\mathcal{U}_S(\Skin)$.
For any history $h=(u,\mathcal{O})\in\mathcal{H}_S$,
\[
\mathrm{Pred}_{\Skin}(h;u_0(\mathbf{p}))
=\mathrm{Pred}_{\Skin}(\mathcal{O};u(u_0(\mathbf{p})))
=\mathrm{Pred}_{\Skin}(\mathcal{O};(u\circ u_0)(\mathbf{p})),
\]
and similarly for $\mathbf{p}'$.
Since $\mathcal{U}_S(\Skin)$ is a monoid, $u\circ u_0\in\mathcal{U}_S(\Skin)$, so the equalities encoded in
$\Sigma_{\Skin}(\mathbf{p})=\Sigma_{\Skin}(\mathbf{p}')$ imply
$\mathrm{Pred}_{\Skin}(h;u_0(\mathbf{p}))=\mathrm{Pred}_{\Skin}(h;u_0(\mathbf{p}'))$ for all $h\in\mathcal{H}_S$.
Thus $\Sigma_{\Skin}(u_0(\mathbf{p}))=\Sigma_{\Skin}(u_0(\mathbf{p}'))$, i.e.\ $u_0(\mathbf{p})\sim_{\ker}u_0(\mathbf{p}')$.

\textbf{(ii)} Let $\approx$ satisfy (a)--(b) and assume $\mathbf{p}\approx \mathbf{p}'$.
Then by (b), $u(\mathbf{p})\approx u(\mathbf{p}')$ for all $u\in\mathcal{U}_S(\Skin)$.
By (a), this implies $u(\mathbf{p})\sim_{0}u(\mathbf{p}')$ for all $u$, i.e.\ for all $\mathcal{O}\in S_0$,
\[
\mathrm{Pred}_{\Skin}(\mathcal{O};u(\mathbf{p}))=\mathrm{Pred}_{\Skin}(\mathcal{O};u(\mathbf{p}')).
\]
By Definition~\ref{def:histories} this is exactly $\Sigma_{\Skin}(\mathbf{p})=\Sigma_{\Skin}(\mathbf{p}')$, hence $\mathbf{p}\sim_{\ker}\mathbf{p}'$.

\textbf{(iii)} If $\mathbf{p}\sim_{0}\mathbf{p}'$ but not $\mathbf{p}\sim_{\ker}\mathbf{p}'$, then by definition there exist
$u\in\mathcal{U}_S(\Skin)$ and $\mathcal{O}\in S_0$ with
$\mathrm{Pred}_{\Skin}(\mathcal{O};u(\mathbf{p}))\neq \mathrm{Pred}_{\Skin}(\mathcal{O};u(\mathbf{p}'))$.
Thus $u(\mathbf{p})\not\sim_{0}u(\mathbf{p}')$, so the induced map $[\mathbf{p}]_{\sim_0}\mapsto [u(\mathbf{p})]_{\sim_0}$ on
$P(\Skin)/\!\sim_0$ is not well-defined: the value depends on the chosen representative.
Treating this dependence as standing-relevant is representational privilege, and attempting to ``repair'' it by retroactive representative selection violates (MC3) irreversible commitment.
Therefore update-compatible identity must refine to $\sim_{\ker}$, which in a history-saturated scope agrees with $\sim_S$.
\end{proof}


\begin{definition}[Continuation-separability condition \textnormal{(CR-Sep)}]\label{def:cr-sep}
Let $(X,r,\mathrm{Pred}_X)$ be a standing presentation for the history-saturated scope determined by
Definition~\ref{def:histories}. Concretely: for every $h\in\mathcal{H}_S$ there is a prediction map
$\mathrm{Pred}_X(h;\cdot):X\to\mathbb{R}$ such that
\[
\mathrm{Pred}_{\Skin}(h;\mathbf{p})=\mathrm{Pred}_X(h;r(\mathbf{p}))
\qquad (\forall\,\mathbf{p}\in P(\Skin)).
\]
Define the \emph{continuation signature} of a standing state $x\in X$ by
\[
\Sigma_X(x) := \big(\mathrm{Pred}_X(h;x)\big)_{h\in\mathcal{H}_S}\ \in\ \mathbb{R}^{\mathcal{H}_S}.
\]
We say the presentation satisfies \textbf{CR-Sep} (continuation-separability) if $\Sigma_X$ is injective:
for any $x\neq x'$ there exists $h\in\mathcal{H}_S$ with $\mathrm{Pred}_X(h;x)\neq \mathrm{Pred}_X(h;x')$.
\end{definition}

\begin{definition}[Continuation-kernel reduction of a standing presentation]\label{def:cr-reduction}
Let $(X,r,\mathrm{Pred}_X)$ be a standing presentation for the history-saturated scope determined by
Definition~\ref{def:histories}. Define an equivalence relation $\approx_S$ on $X$ by
\[
x\approx_S x' \quad\Longleftrightarrow\quad \mathrm{Pred}_X(h;x)=\mathrm{Pred}_X(h;x')\ \text{for all }h\in\mathcal{H}_S .
\]
Let $\pi_X:X\to X^{\mathrm{red}}:=X/\!\approx_S$ be the quotient map. Define the \emph{reduced} representation
\[
r^{\mathrm{red}}:=\pi_X\circ r : P(\Skin)\to X^{\mathrm{red}},
\]
and define reduced predictions by
\begin{equation}\label{eq:PredXred}
\mathrm{Pred}_{X^{\mathrm{red}}}(h;[x])\ :=\ \mathrm{Pred}_X(h;x)\qquad (h\in\mathcal{H}_S).
\end{equation}
\end{definition}

\begin{lemma}[Reduction yields continuation-separability and forces soundness]\label{lem:cr-reduction}
Let $(X,r,\mathrm{Pred}_X)$ be any standing presentation for the history-saturated scope determined by
Definition~\ref{def:histories}, and let $(X^{\mathrm{red}},r^{\mathrm{red}},\mathrm{Pred}_{X^{\mathrm{red}}})$
be its reduction (Definition~\ref{def:cr-reduction}). Then:
\begin{enumerate}[label=(\roman*),leftmargin=2.2em]
\item $(X^{\mathrm{red}},r^{\mathrm{red}},\mathrm{Pred}_{X^{\mathrm{red}}})$ is a standing presentation for $S$.
\item The reduced presentation is sound and complete.
\item The reduced presentation satisfies \textbf{CR-Sep} (Definition~\ref{def:cr-sep}).
\end{enumerate}
Consequently, \textbf{CR-Sep} can be taken \emph{without loss of generality}: any same-scope standing presentation
admits a canonical reduction whose standing states are continuation-separated.
\end{lemma}

\begin{proof}
(i) First note that $\approx_S$ is an equivalence relation (reflexive, symmetric, transitive) by the
corresponding properties of equality of real numbers, quantified over $h\in\mathcal{H}_S$.
To see that \eqref{eq:PredXred} is well-defined, suppose $[x]=[x']$ in $X^{\mathrm{red}}$, i.e.\ $x\approx_S x'$.
Then by definition $\mathrm{Pred}_X(h;x)=\mathrm{Pred}_X(h;x')$ for all $h$, so the right-hand side of
\eqref{eq:PredXred} does not depend on the chosen representative. Surjectivity of $r^{\mathrm{red}}$ follows
because $r$ and $\pi_X$ are surjective. Finally, for any $\mathbf{p}\in P(\Skin)$ and any $h\in\mathcal{H}_S$,
\[
\mathrm{Pred}_{\Skin}(h;\mathbf{p})
=\mathrm{Pred}_X(h;r(\mathbf{p}))
=\mathrm{Pred}_{X^{\mathrm{red}}}\big(h;\pi_X(r(\mathbf{p}))\big)
=\mathrm{Pred}_{X^{\mathrm{red}}}\big(h;r^{\mathrm{red}}(\mathbf{p})\big),
\]
so the factorization condition holds and $(X^{\mathrm{red}},r^{\mathrm{red}},\mathrm{Pred}_{X^{\mathrm{red}}})$
is a standing presentation.

(ii) Soundness: assume $\mathbf{p}\sim_S \mathbf{p}'$. By Definition~\ref{def:standing-equiv} (read as quantifying
over histories via Definition~\ref{def:histories}), we have
$\mathrm{Pred}_{\Skin}(h;\mathbf{p})=\mathrm{Pred}_{\Skin}(h;\mathbf{p}')$ for all $h\in\mathcal{H}_S$.
By factorization this implies
$\mathrm{Pred}_X(h;r(\mathbf{p}))=\mathrm{Pred}_X(h;r(\mathbf{p}'))$ for all $h$, hence
$r(\mathbf{p})\approx_S r(\mathbf{p}')$ and therefore
$r^{\mathrm{red}}(\mathbf{p})=\pi_X(r(\mathbf{p}))=\pi_X(r(\mathbf{p}'))=r^{\mathrm{red}}(\mathbf{p}')$.

Completeness: assume $r^{\mathrm{red}}(\mathbf{p})=r^{\mathrm{red}}(\mathbf{p}')$. Then
$\pi_X(r(\mathbf{p}))=\pi_X(r(\mathbf{p}'))$, i.e.\ $r(\mathbf{p})\approx_S r(\mathbf{p}')$, so
$\mathrm{Pred}_X(h;r(\mathbf{p}))=\mathrm{Pred}_X(h;r(\mathbf{p}'))$ for all $h$ and hence
$\mathrm{Pred}_{\Skin}(h;\mathbf{p})=\mathrm{Pred}_{\Skin}(h;\mathbf{p}')$ for all $h$ by factorization.
Thus $\mathbf{p}\sim_S \mathbf{p}'$.

(iii) Define the continuation signature on $X^{\mathrm{red}}$ by
\[
\Sigma_{X^{\mathrm{red}}}([x]) := \big(\mathrm{Pred}_{X^{\mathrm{red}}}(h;[x])\big)_{h\in\mathcal{H}_S}
=\big(\mathrm{Pred}_X(h;x)\big)_{h\in\mathcal{H}_S}.
\]
If $\Sigma_{X^{\mathrm{red}}}([x])=\Sigma_{X^{\mathrm{red}}}([x'])$, then
$\mathrm{Pred}_X(h;x)=\mathrm{Pred}_X(h;x')$ for all $h$, i.e.\ $x\approx_S x'$, hence $[x]=[x']$.
So $\Sigma_{X^{\mathrm{red}}}$ is injective and the reduced presentation satisfies CR-Sep.
\end{proof}


\begin{lemma}[CR-Sep implies soundness without invoking the weakest reading of MC1]\label{lem:cr-sound}
Let $(X,r,\mathrm{Pred}_X)$ be a standing presentation for scope $S$ and assume \textbf{CR-Sep}
(Definition~\ref{def:cr-sep}).
Then the presentation is sound:
\[
\mathbf{p}\sim_S\mathbf{p}'\ \Longrightarrow\ r(\mathbf{p})=r(\mathbf{p}').
\]
\end{lemma}

\begin{proof}
Assume $\mathbf{p}\sim_S\mathbf{p}'$.
By Definition~\ref{def:standing-presentation}, for every $h\in\mathcal{H}_S$,
\[
\mathrm{Pred}_{\Skin}(h;\mathbf{p})
=\mathrm{Pred}_X(h;r(\mathbf{p})),\qquad
\mathrm{Pred}_{\Skin}(h;\mathbf{p}')
=\mathrm{Pred}_X(h;r(\mathbf{p}')).
\]
Since $\mathbf{p}\sim_S\mathbf{p}'$ implies $\mathrm{Pred}_{\Skin}(h;\mathbf{p})=\mathrm{Pred}_{\Skin}(h;\mathbf{p}')$
for all $h$ (Definition~\ref{def:histories}), we have
$\mathrm{Pred}_X(h;r(\mathbf{p}))=\mathrm{Pred}_X(h;r(\mathbf{p}'))$ for all $h\in\mathcal{H}_S$, i.e.
$\Sigma_X(r(\mathbf{p}))=\Sigma_X(r(\mathbf{p}'))$.
Injectivity of $\Sigma_X$ (CR-Sep) gives $r(\mathbf{p})=r(\mathbf{p}')$.
\end{proof}

\begin{theorem}[Coherent realist identity collapses to the standing quotient]\label{thm:realist-collapse}
Fix a scope declaration $S=(S_0,\mathcal{U}_S)$ and a skin $\Skin$.
Let $(X,r,\mathrm{Pred}_X)$ be any standing presentation for the history-saturated scope determined by
Definition~\ref{def:histories}, and let $(X^{\mathrm{red}},r^{\mathrm{red}},\mathrm{Pred}_{X^{\mathrm{red}}})$
be its reduction (Definition~\ref{def:cr-reduction}).
Then there exists a unique bijection
\[
\phi:\ P(\Skin)/\!\sim_S\ \longrightarrow\ X^{\mathrm{red}},\qquad \phi([\mathbf{p}]):=r^{\mathrm{red}}(\mathbf{p}).
\]
In particular, the identity-bearing content of any same-scope standing presentation is (up to unique isomorphism)
the standing quotient $P(\Skin)/\!\sim_S$; any additional same-scope structure in $X$ beyond its reduction
$X^{\mathrm{red}}$ is continuation-inert label redundancy.
\end{theorem}

\begin{proof}
By Lemma~\ref{lem:cr-reduction}, the reduced representation $r^{\mathrm{red}}$ is sound and complete.

Well-definedness of $\phi$: if $[\mathbf{p}]=[\mathbf{p}']$ in $P(\Skin)/\!\sim_S$, then $\mathbf{p}\sim_S\mathbf{p}'$
and soundness gives $r^{\mathrm{red}}(\mathbf{p})=r^{\mathrm{red}}(\mathbf{p}')$.

Injectivity: if $\phi([\mathbf{p}])=\phi([\mathbf{p}'])$, then
$r^{\mathrm{red}}(\mathbf{p})=r^{\mathrm{red}}(\mathbf{p}')$ and completeness implies $\mathbf{p}\sim_S\mathbf{p}'$,
hence $[\mathbf{p}]=[\mathbf{p}']$.

Surjectivity: $r^{\mathrm{red}}$ is surjective because $r$ and $\pi_X$ are surjective.

Uniqueness: if $\psi:P(\Skin)/\!\sim_S\to X^{\mathrm{red}}$ also satisfies $\psi([\mathbf{p}])=r^{\mathrm{red}}(\mathbf{p})$,
then $\psi=\phi$ because the quotient map $q:P(\Skin)\to P(\Skin)/\!\sim_S$ is surjective
(Lemma~\ref{lem:quotient-universal}).
\end{proof}


\begin{remark}[Relation to $\Ad_S$ and $\mathrm{NF}_S$]\label{rem:cr-nf}
Theorem~\ref{thm:realist-collapse} shows that a realist can accept additional \emph{unmodeled} structure only by
either enlarging scope (admitting new histories/observables) or adding new load (new witnesses); otherwise any
same-scope identity refinement is eliminated by the reduction of Definition~\ref{def:cr-reduction} (Lemma~\ref{lem:cr-reduction}) and becomes haecceitistic.
In particular, duality-related reorganizations that preserve all $S$-predictions cannot introduce new standing
witnesses beyond the canonical normal form $\mathrm{NF}_S$ (Definition~\ref{def:NF-skin}) and correspond to
isomorphisms in the admissibility groupoid $\Ad_S$ (Theorem~\ref{thm:coherence-groupoid}).
\end{remark}



\begin{corollary}[Transport completeness is forced by ``same-scope'' claims]\label{cor:transport-forced}
Let $\Skin_A,\Skin_B$ be skins both used for scope $S$.
If one asserts that $\Skin_A$ and $\Skin_B$ represent the \emph{same} standing content for $S$
(i.e.\ there exists a presentation $X$ for which both factor as in \eqref{eq:standing-presentation}),
then any admissible transport implementing that assertion must induce a bijection on standing quotients
and is therefore transport-complete for $S$ (Definition~\ref{def:transport-complete}).
Equivalently: within minimal coherence, ``empirical equivalence on $S$'' plus ``same-scope identity''
already implies transport completeness; it is not an optional strengthening.
\end{corollary}

\begin{proof}
By Theorem~\ref{thm:standing-quotient-unique}, each skin's standing state space is uniquely
$P(\Skin_\bullet)/\!\sim_S$ up to bijection.
A transport implementing ``same-scope identity'' must therefore identify standing classes bijectively.
If it failed transport completeness, some $S$-distinguishable pair in one skin would map to an
$S$-indistinguishable pair in the other, contradicting the uniqueness of standing state spaces and violating
stable reference (MC1/MC3). Hence any same-scope transport must be transport-complete.
\end{proof}

\begin{theorem}[No nontrivial standing-repair or canonical-selection operators]\label{thm:no-repair}
Fix a scope $S$ and assume minimal coherence.
There is no admissible operator that, without enlarging the scope or adding new load, refines standing beyond
$\sim_S$ or selects a standing-relevant canonical representative from each standing class.
More precisely:
\begin{enumerate}[label=(\roman*),leftmargin=2.2em]
\item Any rule that distinguishes two points $\mathbf{p}\sim_S\mathbf{p}'$ as different identity-bearing
states is a standing failure (Definition~\ref{def:failure}).
\item Any ``canonical representative'' selector $s:P(\Skin)/\!\sim_S\to P(\Skin)$ (a section of the quotient map)
may be used as a calculational convenience (gauge choice), but cannot be standing-relevant:
treating $s([\mathbf{p}])$ as identity-grounding constitutes representational privilege and breaks (MC1).
\item Consequently, typicality/measure/choice-selection operators that require a canonical representative or a
non-invariant measure on representatives are inadmissible inside the same-scope interior; using them is an
illicit scope transport (Definition~\ref{def:failure}(c)).
\end{enumerate}
\end{theorem}

\begin{proof}
(i) If $\mathbf{p}\sim_S\mathbf{p}'$ then by definition no $S$-observable witnesses any distinction between them.
Declaring them different identity-bearing states therefore introduces hidden unwitnessed structure or privileged
representation, i.e.\ a standing failure.

(ii) A section $s$ chooses one representative per standing class. If $s$ is treated as standing-relevant, then the
identity of the referent depends on the choice of section, i.e.\ on a privileged representation. This violates
(MC1). Therefore $s$ can exist only as gauge convenience: any two choices of section are admissibly related and
must not be used to generate new standing content.

(iii) A typicality or choice-selection operator requires additional structure beyond the standing class to define
``most'' or ``selected'' representatives (a measure, reference class, selector).
If that structure is not invariant under admissible redescription, it is representational privilege; if it is
kept fixed anyway, it imports non-witnessed anchoring and is therefore an illicit scope transport.
Either way it is inadmissible inside the same-scope interior.
\end{proof}

\begin{remark}[Common ``standing repairs'' and why they fail fast]\label{rem:repair-catalog}
Hostile reports often propose a meta-level ``repair'' to evade a closure result: add a selector, add a measure,
add a completion operator, add a genericity/typicality rule, etc.
Inside the present governance discipline these all reduce to the same typed failure:
they attempt to introduce standing-relevant structure not fixed by $\sim_S$ and therefore not witnessed in scope $S$.
Typical examples include:
\begin{itemize}[leftmargin=2.2em]
\item \emph{Canonical representative selection} (choose a preferred element of each $[\mathbf{p}]$): violates
Theorem~\ref{thm:no-repair}(ii) unless treated as pure gauge convenience.
\item \emph{Typicality/measure postulates} on representatives (``most'' points in a class): require a non-invariant
measure or reference-class choice and are therefore representational privilege or illicit scope transport.
\item \emph{Choice-selection operators} (pick one completion/branch when multiple are $S$-indistinguishable):
introduce hidden unwitnessed structure unless the choice is itself witnessed (scope change).
\item \emph{Completion/totalization operators} that claim to make partial data ``total'' without new load:
they are either redundant relabelings of $\mathrm{NF}_S$ (no new standing) or alter witness content (scope change).
\item \emph{Genericity forcing / saturation / compactness-as-operator} repairs: these are meta-existence moves that
select a privileged enlargement or completion; absent new standing witnesses they collapse back to representational
privilege, and with witnesses they are scope changes.
\end{itemize}
The paper therefore fail-closes: such repairs are admissible only if they are exhibited as explicit load and scope
changes with new witnesses, not as same-scope ``equivalence completions.''
\end{remark}




\begin{corollary}[Duality reduction]\label{cor:duality-reduction}
Let $\Skin_A,\Skin_B$ be any two skins (including duality-related reorganizations) that agree on all
$S$-predictions and preserve standing under minimal coherence.
Then they determine the same standing quotient up to isomorphism:
\[
P(\Skin_A)/\!\sim_S\ \cong\ P(\Skin_B)/\!\sim_S,
\]
and hence cannot introduce additional invariant witnesses for $S$ beyond those already present in the
canonical normal form $\mathrm{NF}_S$ (Definition~\ref{def:NF-skin}).
\end{corollary}

\begin{proof}
By Theorem~\ref{thm:standing-quotient-unique}, each skin's standing state space is uniquely its quotient.
Agreement on all $S$-predictions implies the quotients have identical prediction rule, hence are isomorphic as
standing presentations. Any purported new witness would refine standing beyond $\sim_S$, contradicting
Theorem~\ref{thm:no-repair}(i).
\end{proof}

\begin{theorem}[Maximality: no proper standing-preserving extension of $\Ad_S$]\label{thm:maximality}
Fix a scope $S$ and assume minimal coherence.
Let $\Ad_S$ be the admissibility groupoid of same-scope equivalences (Theorem~\ref{thm:coherence-groupoid}).
There is no proper enlargement of $\Ad_S$ that preserves the same scope $S$, introduces no new load,
and remains minimally coherent. Any additional ``operator'' or ``morphism'' either:
\begin{enumerate}[label=(\roman*),leftmargin=2.2em]
\item factors through the standing quotient and therefore adds no new standing content (redundancy), or
\item changes the witness content and is therefore a scope change, or
\item violates minimal coherence via a standing-failure mode.
\end{enumerate}
\end{theorem}

\begin{proof}
Let $F$ be any additional same-scope map introduced on top of $\Ad_S$.
If $F$ preserves $S$-predictions, then it cannot distinguish $\sim_S$-equivalent points: if
$\mathbf{p}\sim_S\mathbf{p}'$ but $F(\mathbf{p})\neq F(\mathbf{p}')$, then $F$ would introduce an identity-bearing
distinction with no $S$-witness, contradicting Theorem~\ref{thm:no-repair}(i). Thus $F$ is constant on
$\sim_S$-classes. Let $q:P(\Skin)\to P(\Skin)/\!\sim_S$ be the quotient map and define
$\widetilde F([\mathbf{p}]):=F(\mathbf{p})$. This is well-defined and satisfies $F=\widetilde F\circ q$; uniqueness
follows from surjectivity of $q$. Hence $F$ adds no standing content beyond the quotient and is redundancy.

If $F$ does not preserve $S$-predictions, then by definition it changes the witness content and is a scope change.
If $F$ purports to preserve $S$ while distinguishing $\sim_S$-equivalent points or selecting standing-relevant
representatives, it violates Theorem~\ref{thm:no-repair} and hence minimal coherence.
Therefore there is no proper same-scope standing-preserving extension of $\Ad_S$.
\end{proof}

\subsection{Typed primitives}

\begin{definition}[Bookkeeping skin]\label{def:skin}
A \emph{bookkeeping skin} $\Skin$ is a regime-scoped representational scheme: it specifies
(i) a choice of variables/fields/parameters, and (ii) rules for computing predictions or invariant
quantities in that regime.
\end{definition}



\begin{definition}[Scope]\label{def:scope}
A \emph{scope} $S$ is a specification of:
\begin{enumerate}[label=(\roman*),leftmargin=2.2em]
\item a regime (energy/momentum range, on-shell/off-shell conditions, and degrees retained);
\item a class of admissible questions/observables (e.g.\ low-energy CC contact amplitudes; pure QED
scattering; on-shell high-energy longitudinal amplitudes);
\item an admissible error discipline, if approximations are used (e.g.\ ``up to $\cO(p^2/m_W^2)$'').
\end{enumerate}
All standing and witness claims in this paper are \emph{scope-indexed}: a parameter can be redundant for
one scope and standing-relevant for a larger scope.
\end{definition}

\begin{definition}[Scope-equivalent skins]\label{def:scope-equivalent}
Two skins $\Skin_A,\Skin_B$ are \emph{scope-equivalent} for $S$ if there exist admissible transports
$T_{A\to B}$ and $T_{B\to A}$ such that, for every observable $\mathcal{O}$ admitted by $S$,
the predictions agree:
\[
\mathrm{Pred}_{A}(\mathcal{O};\mathbf{p})=\mathrm{Pred}_{B}(\mathcal{O};T_{A\to B}(\mathbf{p}))
\quad\text{and}\quad
\mathrm{Pred}_{B}(\mathcal{O};\mathbf{q})=\mathrm{Pred}_{A}(\mathcal{O};T_{B\to A}(\mathbf{q})).
\]
\end{definition}

\begin{definition}[Transport completeness]\label{def:transport-complete}
An admissible transport $T:\Skin_A\to\Skin_B$ is \emph{transport-complete} for scope $S$ if $\Skin_B$
retains enough data to reproduce \emph{all} standing-relevant distinctions of $\Skin_A$ that affect
$S$-observables. Equivalently: whenever two parameter points $\mathbf{p},\mathbf{p}'$ in $\Skin_A$ yield
different predictions for some $\mathcal{O}\in S$, their transports $T(\mathbf{p}),T(\mathbf{p}')$ differ
in at least one invariant witness of $\Skin_B$.
\end{definition}

\begin{remark}[EFT truncation is not ``admissible equivalence'' without an error envelope]\label{rem:eft}
An $\EFT$ may be an admissible skin for a scope $S$ only if it is \emph{matching-complete for $S$}:
all operators whose coefficients affect $S$ within the declared error tolerance are included, or else an
explicit error envelope is carried as part of the witness bookkeeping.
A truncated $\EFT$ presented as exact equivalence is inadmissible under the fail-closed discipline.
\end{remark}

\begin{remark}[No existence-by-assertion of admissible skins]\label{rem:no-existence-by-assertion}
Admissibility is fail-closed: the existence of an admissible skin, transport, or ``further equivalence'' is
not a matter of plausibility. Any alleged additional admissible move must be exhibited by explicit construction
and supplied with a transport certificate (Definition~\ref{def:certificate}) inside the declared load and the
declared scope. Bare existential claims (``there may exist some other admissible skin not considered'') do not
reopen the space.
\end{remark}


\begin{definition}[Standing]\label{def:standing}
\emph{Standing} is the referential identity of the quantities tracked by the bookkeeping.
A datum has standing if it is invariant under all admissible redescriptions \emph{and} is transportable
across admissible regime interfaces without importing extra load.
\end{definition}

\begin{definition}[Admissible redescription]\label{def:redescription}
An \emph{admissible redescription} within a skin is any transformation that preserves (i) regime of
applicability and (ii) standing of tracked referents. Examples in \SM bookkeeping include:
field redefinitions, unitary flavor-basis rotations, rephasings, gauge-fixing changes, and
reparameterizations that leave invariant content unchanged.
\end{definition}

\begin{definition}[Admissible transport]\label{def:transport}
A \emph{transport map} $T:\Skin_A\to\Skin_B$ is an explicit rule mapping data in $\Skin_A$ to data in
$\Skin_B$ that describe the same physical content in an overlapping regime (matching, integrating out,
unbroken$\leftrightarrow$broken transport, on-shell reduction, etc.).
Transport is \emph{admissible} if it preserves standing and introduces no new load outside
Definition~\ref{def:declared-load}.
\end{definition}

\begin{definition}[Invariant witness set]\label{def:witness}
Let $\Skin$ be a skin. An \emph{invariant witness set} $\Wit(\Skin)$ is a set of quantities
(or equivalence classes) such that:
\begin{enumerate}[label=(\roman*),leftmargin=2.2em]
\item each element of $\Wit(\Skin)$ is invariant under every admissible redescription in $\Skin$;
\item for every admissible transport $T:\Skin\to\Skin'$, the transported standing-relevant content depends
only on the images of $\Wit(\Skin)$.
\end{enumerate}
\end{definition}

\begin{definition}[Witness rank]\label{def:witness-rank}
Fix a scope $S$.
The \emph{witness rank} for $S$ is the minimal \emph{cardinality} (finite or infinite) of independent real
parameters required to specify a minimal invariant witness set sufficient to determine all $S$-predictions,
with discrete skeleton labels held fixed.
Independence is understood within the explicit certificate discipline of this paper: admissible witnesses and
transports must be finitely presented and evaluable from the declared load, so pathological ``oracle encodings''
(e.g.\ hiding unbounded information in digits of a single real) are excluded as inadmissible.
(Throughout, ``rank'' comparisons concern this continuous coordinate cardinality; discrete skeleton data are
treated separately.)
\end{definition}

\begin{remark}[Finite vs.\ functional witness rank]\label{rem:witness-rank-functional}
The witness rank may be infinite when the admitted scope requires genuinely functional data (e.g.\ a spectral
density rather than finitely many numbers). In such cases, ``rank'' comparisons are understood as cardinal
comparisons. See Proposition~\ref{prop:em-hvp} for an explicit example in the electromagnetic sector.
\end{remark}

\begin{remark}[Finite-dimensional regular case]\label{rem:witness-rank-regular}
In the finite-parameter sectors treated in this paper, prediction maps and witness functions are analytic in the
continuous parameters away from thresholds. In that regular setting the witness-rank bound of
Theorem~\ref{thm:witness-counting} matches ordinary parameter counting: if the standing quotient carries the structure
of a finite-dimensional smooth manifold, any smooth injective witness map into $\mathbb{R}^m$ implies
$\dim\!\left(P(\Skin)/\!\sim_S\right)\le m$ (invariance of domain / rank theorem). The cardinal formulation is retained
only to accommodate scopes with functional witnesses (Remark~\ref{rem:witness-rank-functional}).
\end{remark}


\begin{definition}[Standing equivalence for a scope]\label{def:standing-equiv}
Fix a scope $S$ and a skin $\Skin$ with parameter space $P(\Skin)$.
For $\mathbf{p},\mathbf{p}'\in P(\Skin)$ write $\mathbf{p}\sim_S \mathbf{p}'$ iff, for every observable
$\mathcal{O}$ admitted by $S$,
\[
\mathrm{Pred}_{\Skin}(\mathcal{O};\mathbf{p})=\mathrm{Pred}_{\Skin}(\mathcal{O};\mathbf{p}')
\]
within the error discipline declared by $S$ (Definition~\ref{def:scope}).
\end{definition}

\begin{definition}[Standing state space / quotient]\label{def:standing-quotient}
The \emph{standing state space} of $\Skin$ for scope $S$ is the quotient set $P(\Skin)/\!\sim_S$.
No canonical representative is assumed or required: only equivalence classes are used.
\end{definition}

\begin{lemma}[Scope-equivalence induces an isomorphism of standing quotients]\label{lem:quotient-iso}
Let $\Skin_A,\Skin_B$ be scope-equivalent skins for $S$ (Definition~\ref{def:scope-equivalent}).
Assume there are admissible transports $T_{A\to B}$ and $T_{B\to A}$ that are transport-complete for $S$
(Definition~\ref{def:transport-complete}).
Then $T_{A\to B}$ induces a bijection of standing state spaces
\[
P(\Skin_A)/\!\sim_S\;\cong\; P(\Skin_B)/\!\sim_S.
\]
\end{lemma}

\begin{proof}
If $\mathbf{p}\sim_S \mathbf{p}'$ in $\Skin_A$, then for every $\mathcal{O}\in S$ we have
$\mathrm{Pred}_A(\mathcal{O};\mathbf{p})=\mathrm{Pred}_A(\mathcal{O};\mathbf{p}')$.
By scope-equivalence, $\mathrm{Pred}_B(\mathcal{O};T_{A\to B}(\mathbf{p}))=
\mathrm{Pred}_A(\mathcal{O};\mathbf{p})$ and similarly for $\mathbf{p}'$, hence
$T_{A\to B}(\mathbf{p})\sim_S T_{A\to B}(\mathbf{p}')$.
Thus $T_{A\to B}$ is well-defined on the quotient.

If $\mathbf{p}\not\sim_S \mathbf{p}'$, then some $\mathcal{O}\in S$ differs in prediction.
Transport completeness implies the transported points differ in standing-relevant witness content for $S$ in
$\Skin_B$, hence are not $\sim_S$-equivalent. Therefore the induced map on quotients is injective.
Applying the same argument to $T_{B\to A}$ yields an inverse on quotients, hence bijectivity.
\end{proof}



\begin{definition}[Tensor load]\label{def:tensor-load}
For a fixed scope, the \emph{tensor load} is the minimal data required to fix all invariant witnesses
across all admissible skins for that scope. Parameters not required to fix witnesses have no standing
for that scope (redundancy) unless supported by additional invariant witnesses \emph{in every} admissible skin.
\end{definition}

\subsection{Standing failures: the diagnostic taxonomy}

\begin{definition}[Standing-failure modes]\label{def:failure}
A purported independent degree of freedom is a \emph{standing failure} if maintaining its independence
requires at least one of:
\begin{enumerate}[label=(\alph*),leftmargin=2.2em]
\item \textbf{hidden unwitnessed structure}: dependence on data not represented by any invariant witness;
\item \textbf{representational privilege}: privileging a particular skin/basis/parameterization as identity-grounding;
\item \textbf{illicit scope transport}: importing structure justified only outside the declared scope;
\item \textbf{compensating load/patching}: adding fields/operators/sectors forbidden by the declared load.
\end{enumerate}
\end{definition}

\subsection{The engine theorem}

\begin{theorem}[Standing rank bound (witness counting)]\label{thm:witness-counting}
Fix a scope $S$ (Definition~\ref{def:scope}) and the declared load.
Let $\Skin_A,\Skin_B$ be scope-equivalent skins for $S$ (Definition~\ref{def:scope-equivalent}) and assume
there exist admissible transports between them that are transport-complete for $S$
(Definition~\ref{def:transport-complete}).
Let $P(\Skin_A)$ be parameterized by $n$ real coordinates $\mathbf{p}=(p_1,\dots,p_n)$, with all discrete
skeleton labels held fixed.

Let $m$ be the witness rank for $S$ (Definition~\ref{def:witness-rank}); equivalently, $m$ is the minimal
cardinality of real parameters needed to coordinatize the standing state space $P(\Skin_B)/\!\sim_S$
(Definition~\ref{def:standing-quotient}).
Then $\Skin_A$ cannot carry more than $m$ independent standing-relevant continuous degrees of freedom for $S$.
In particular, if $n>m$ then $\mathbf{p}$ contains redundancy directions for $S$; maintaining them as
``standing'' requires a standing failure (Definition~\ref{def:failure}).
\end{theorem}

\begin{proof}
By Lemma~\ref{lem:quotient-iso}, scope-equivalence with transport completeness induces a bijection
\[
P(\Skin_A)/\!\sim_S\ \cong\ P(\Skin_B)/\!\sim_S.
\]
Hence it suffices to bound the standing coordinate cardinality in $\Skin_B$.

Let $\Wit_S(\Skin_B)=\{w_i\}_{i\in I}$ be a \emph{minimal} invariant witness set determining all $S$-predictions,
where $I$ is an index set of cardinality $|I|=m$ (Definition~\ref{def:witness-rank}).
Define the \emph{witness evaluation map}
\[
w_B:\ P(\Skin_B)\longrightarrow \mathbb{R}^{I},\qquad
w_B(\mathbf{p}) := \big(w_i(\mathbf{p})\big)_{i\in I},
\]
where $\mathbb{R}^{I}$ denotes the set of functions $I\to\mathbb{R}$ (finite-dimensional vectors when $I$ is finite).

Because $\Wit_S(\Skin_B)$ determines all $S$-predictions, for each $\mathcal{O}\in S$ there exists a function
$F_{\mathcal{O}}:\mathbb{R}^{I}\to\mathbb{R}$ such that
\[
\mathrm{Pred}_{\Skin_B}(\mathcal{O};\mathbf{p}) \;=\; F_{\mathcal{O}}\!\big(w_B(\mathbf{p})\big)
\qquad (\forall\,\mathbf{p}\in P(\Skin_B)).
\]
In particular, if $w_B(\mathbf{p})=w_B(\mathbf{p}')$, then $\mathrm{Pred}_{\Skin_B}(\mathcal{O};\mathbf{p})
=\mathrm{Pred}_{\Skin_B}(\mathcal{O};\mathbf{p}')$ for all $\mathcal{O}\in S$, hence
$\mathbf{p}\sim_S\mathbf{p}'$ by Definition~\ref{def:standing-equiv}.
Thus $w_B$ is constant on $\sim_S$-classes and induces a well-defined map on the quotient,
\[
\overline w_B:\ P(\Skin_B)/\!\sim_S\ \longrightarrow\ \mathbb{R}^{I},
\qquad \overline w_B([\mathbf{p}]):=w_B(\mathbf{p}).
\]

Moreover, $\overline w_B$ is injective: if $[\mathbf{p}]\neq[\mathbf{p}']$ then $\mathbf{p}\not\sim_S\mathbf{p}'$,
so some $\mathcal{O}\in S$ has $\mathrm{Pred}_{\Skin_B}(\mathcal{O};\mathbf{p})\neq
\mathrm{Pred}_{\Skin_B}(\mathcal{O};\mathbf{p}')$.
But then necessarily $w_B(\mathbf{p})\neq w_B(\mathbf{p}')$, otherwise all $S$-predictions would coincide by the
display above. Hence $\overline w_B([\mathbf{p}])\neq \overline w_B([\mathbf{p}'])$.

Therefore $P(\Skin_B)/\!\sim_S$ embeds into $\mathbb{R}^{I}$ and can carry at most $|I|=m$ independent real
coordinates (in the cardinal sense when $m$ is infinite). The same bound holds for $\Skin_A$ because its
standing quotient is bijective to that of $\Skin_B$.

If $\Skin_A$ is written with $n>m$ real coordinates, then the quotient map
$P(\Skin_A)\to P(\Skin_A)/\!\sim_S$ cannot be injective: there exist distinct parameter points in $P(\Skin_A)$
representing the same standing class. Any coordinate direction along such fibers is redundancy for scope $S$.
Maintaining such coordinates as ``standing'' would require either denying transport completeness/scope-equivalence,
or importing additional standing structure (new witnesses, privileged parameterization, or hidden load),
i.e.\ a standing-failure mode (Definition~\ref{def:failure}).
\end{proof}

\begin{remark}[Scope-indexed, not global]\label{rem:scope-index}
Theorem~\ref{thm:witness-counting} does \emph{not} claim that a parameter redundant for one scope is
redundant in all scopes. Enlarging $S$ (e.g.\ by including higher-energy observables) can increase the
minimal witness rank and restore standing to parameters previously redundant in a restricted regime.
\end{remark}


\subsection{Transport certificates: how every collapse will be audited}

\begin{definition}[Transport certificate schema]\label{def:certificate}
A \emph{transport certificate} for a collapse lemma consists of:
\begin{enumerate}[label=(\textbf{C\arabic*}),leftmargin=2.8em]
\item explicit identification of the input skin $\Skin_A$ (variables and parameters);
\item explicit identification of the output skin $\Skin_B$ (witness quantities);
\item an explicit admissible transport $T:\Skin_A\to\Skin_B$ (matching, reduction, or redescription);
\item an explicit computation of $\Wit(\Skin_B)$ and its rank;
\item the rank comparison and the resulting redundancy/standing-failure diagnosis by
Theorem~\ref{thm:witness-counting}.
\end{enumerate}
\end{definition}

\section{Parameter normal form: fully worked witness collapses}

\subsection{Charged-current collapse to $G_F$ (fully worked)}

\begin{lemma}[Charged-current collapse in the contact-interaction scope]\label{lem:cc}
Let $S_{\CC}^{(6)}$ be the \emph{charged-current contact scope} defined by $p^2\ll m_W^2$ and retention
of the leading dimension-6 four-fermion operator, with corrections treated as
$\cO(p^2/m_W^2)$ (Definition~\ref{def:scope}).
In $S_{\CC}^{(6)}$, admissible transport from the renormalizable mediator skin to the four-fermion
$\EFT$ skin exposes a single invariant witness $G_F$. Consequently, the pair $(g_2,v)$ overcounts
tensor load for \emph{this restricted scope}: only the combination $G_F/\sqrt2=g_2^2/(8m_W^2)$ has
standing as charged-current strength in $S_{\CC}^{(6)}$.
\end{lemma}

\begin{proof}
We present a transport certificate in the sense of Definition~\ref{def:certificate}.

\paragraph{(C1) Input skin $\Skin_W$ (mediator bookkeeping).}
In the electroweak gauge skin, charged currents couple to $W^\pm$ with coupling $g_2$:
\begin{equation}\label{eq:LintW}
\mathcal{L}_{\text{int}}^{\CC} \;=\; -\frac{g_2}{\sqrt{2}}\Big(W^+_\mu J_-^\mu + W^-_\mu J_+^\mu\Big),
\qquad
J_+^\mu := \bar\psi_1\gamma^\mu P_L \psi_2,
\end{equation}
where $P_L=(1-\gamma_5)/2$ and $J_-^\mu=(J_+^\mu)^\dagger$.

\paragraph{(C2) Mediator mass from the Higgs kinetic term.}
Write the Higgs doublet in unitary gauge as
\begin{equation}
H(x)=\frac{1}{\sqrt{2}}
\begin{pmatrix}
0\\ v+h(x)
\end{pmatrix}.
\end{equation}
The gauge-covariant derivative is
\begin{equation}
D_\mu = \partial_\mu - i\frac{g_2}{2}\tau^a W^a_\mu - i g_1 Y B_\mu.
\end{equation}
The Higgs kinetic term $(D_\mu H)^\dagger(D^\mu H)$ contains the charged gauge-boson mass term.
Using $W^\pm_\mu=(W^1_\mu\mp i W^2_\mu)/\sqrt2$ and the Pauli matrices,
one finds the quadratic term
\begin{equation}\label{eq:mWterm}
(D_\mu H)^\dagger(D^\mu H)\supset \frac{g_2^2 v^2}{4}\,W^+_\mu W^{-\mu}.
\end{equation}
Therefore
\begin{equation}\label{eq:mW}
m_W^2=\frac{g_2^2 v^2}{4}, \qquad m_W=\frac{g_2 v}{2}.
\end{equation}

\paragraph{(C3) Output skin $\Skin_F$ (low-energy four-fermion bookkeeping).}
In the low-energy charged-current $\EFT$ skin,
\begin{equation}\label{eq:GFdef}
\mathcal{L}_{\EFT}^{\CC} \;=\; -\frac{G_F}{\sqrt2}\,
\big(\bar\psi_1\gamma^\mu (1-\gamma_5)\psi_2\big)
\big(\bar\psi_3\gamma_\mu (1-\gamma_5)\psi_4\big).
\end{equation}
Here $G_F$ is the \emph{single} coupling witnessing charged-current strength in this regime.

\paragraph{(C4) Admissible transport $T:\Skin_W\to\Skin_F$ (integrate out $W$).}
In the regime $p^2\ll m_W^2$, kinetic terms for $W$ are subleading for matching and we may integrate out
$W$ at tree level by solving its algebraic equation of motion from the mass plus interaction terms:
\begin{equation}
\mathcal{L}\supset m_W^2 W^+_\mu W^{-\mu} -\frac{g_2}{\sqrt2}\Big(W^+_\mu J_-^\mu+W^-_\mu J_+^\mu\Big).
\end{equation}
Varying with respect to $W^-_\mu$ gives
\begin{equation}
m_W^2 W^{+\mu} - \frac{g_2}{\sqrt2}J_+^\mu =0
\quad\Rightarrow\quad
W^{+\mu}=\frac{g_2}{\sqrt2 m_W^2}J_+^\mu,
\end{equation}
and similarly $W^{-\mu}=\frac{g_2}{\sqrt2 m_W^2}J_-^\mu$. Substituting back yields the effective operator
\begin{equation}\label{eq:LeffJJ}
\mathcal{L}_{\mathrm{eff}}^{\CC}
=
-\frac{g_2^2}{2m_W^2}\, J_+^\mu J_{-\mu}.
\end{equation}
Now rewrite the left-handed currents in terms of $(1-\gamma_5)$:
\begin{equation}
J_+^\mu J_{-\mu}
=
\big(\bar\psi_1\gamma^\mu P_L\psi_2\big)\big(\bar\psi_3\gamma_\mu P_L\psi_4\big)
=
\frac{1}{4}\big(\bar\psi_1\gamma^\mu (1-\gamma_5)\psi_2\big)\big(\bar\psi_3\gamma_\mu (1-\gamma_5)\psi_4\big).
\end{equation}
Hence
\begin{equation}\label{eq:matchGF}
\mathcal{L}_{\mathrm{eff}}^{\CC}
=
-\frac{g_2^2}{8m_W^2}\,
\big(\bar\psi_1\gamma^\mu (1-\gamma_5)\psi_2\big)\big(\bar\psi_3\gamma_\mu (1-\gamma_5)\psi_4\big),
\end{equation}
and matching to \eqref{eq:GFdef} gives
\begin{equation}\label{eq:GFmatch}
\frac{G_F}{\sqrt2}=\frac{g_2^2}{8m_W^2}.
\end{equation}
Using \eqref{eq:mW} yields
\begin{equation}\label{eq:GFv}
\frac{G_F}{\sqrt2}=\frac{g_2^2}{8(g_2^2 v^2/4)}=\frac{1}{2v^2}
\quad\Rightarrow\quad
G_F=\frac{1}{\sqrt2\,v^2}.
\end{equation}

\paragraph{(C5) Witness rank and collapse.}
The low-energy $\CC$ skin witnesses interaction strength by the single parameter $G_F$.
The mediator skin uses $(g_2,v)$, but transport collapses all charged-current dependence to the
single combination in \eqref{eq:GFmatch} (equivalently \eqref{eq:GFv}).

Therefore $\Wit(\Skin_F)$ has rank $1$ while $(g_2,v)$ has rank $2$ in the mediator parameterization.
By Theorem~\ref{thm:witness-counting}, $(g_2,v)$ cannot represent two independent standing-relevant
degrees for low-energy charged-current phenomena. The extra degree is redundancy for this scope.
\end{proof}

\begin{remark}[Scope warning for Lemma~\ref{lem:cc}]\label{rem:cc-scope}
Lemma~\ref{lem:cc} is deliberately restricted to the contact-interaction scope $S_{\CC}^{(6)}$.
If the scope is enlarged to resolve momentum dependence of the $W$ propagator (or higher-dimension
operators in the $\EFT$), additional invariant witnesses appear (e.g.\ a witness for $m_W$),
and the ``collapse to $G_F$ alone'' no longer holds. This is an instance of
Remark~\ref{rem:scope-index}.
\end{remark}

\begin{proposition}[Worked example: enlarging $S_{\CC}^{(6)}$ forces an additional witness]\label{prop:cc-p2}
Let $S_{\CC}^{(8)}$ be the \emph{extended charged-current scope} in which one resolves the leading
momentum dependence of the $W$ propagator: one retains the first correction of order $p^2/m_W^2$
beyond the contact limit, with residual errors treated as $\cO(p^4/m_W^4)$.
Equivalently, in the $\EFT$ skin one retains the dimension-6 four-fermion operator together with the
leading derivative-dressed (dimension-8) correction required to reproduce the $p^2/m_W^2$ term.
Then $S_{\CC}^{(8)}$ contains \emph{two} independent invariant witnesses, which may be taken as
$\{G_F,m_W\}$ (or equivalently $\{C_6,C_8\}$, the coefficients of the dimension-6 and dimension-8 operators).
In particular, $m_W$ is an invariant witness in $S_{\CC}^{(8)}$ and cannot be eliminated as redundancy in
that enlarged scope.
\end{proposition}

\begin{proof}
Work at tree level and consider a generic left-handed current--current process mediated by $W$
exchange with momentum transfer $p$.
The full mediator skin gives the amplitude factor (suppressing spinors and group factors already fixed
in Lemma~\ref{lem:cc})
\begin{equation}\label{eq:cc-prop-factor}
\mathcal{M}(p^2)\;\propto\; -\frac{g_2^2}{8}\,\frac{1}{m_W^2-p^2}\,\mathcal{O}_6 \;+\;\cdots ,
\end{equation}
where $\mathcal{O}_6$ denotes the Dirac/field structure
\[
\mathcal{O}_6 :=
\big(\bar\psi_1\gamma^\mu(1-\gamma_5)\psi_2\big)\big(\bar\psi_3\gamma_\mu(1-\gamma_5)\psi_4\big),
\]
and the ellipsis includes terms proportional to $p_\mu p_\nu/m_W^2$ in the massive-vector propagator.
Those terms induce additional derivative operators but are suppressed by the same scale $m_W$; they do
not change the witness-counting point made here.

For $p^2\ll m_W^2$, expand the propagator denominator:
\begin{equation}\label{eq:prop-expand}
\frac{1}{m_W^2-p^2}
=
\frac{1}{m_W^2}\,\frac{1}{1-p^2/m_W^2}
=
\frac{1}{m_W^2}\left(1+\frac{p^2}{m_W^2}+\cO\!\left(\frac{p^4}{m_W^4}\right)\right).
\end{equation}
Substituting \eqref{eq:prop-expand} into \eqref{eq:cc-prop-factor} yields
\begin{equation}\label{eq:cc-expanded}
\mathcal{M}(p^2)\;\propto\;
-\frac{g_2^2}{8m_W^2}\left(1+\frac{p^2}{m_W^2}+\cO\!\left(\frac{p^4}{m_W^4}\right)\right)\mathcal{O}_6.
\end{equation}
Matching the $p^0$ term reproduces Lemma~\ref{lem:cc}:
\begin{equation}\label{eq:C6def}
C_6 := \frac{G_F}{\sqrt2}=\frac{g_2^2}{8m_W^2}.
\end{equation}
In the enlarged scope $S_{\CC}^{(8)}$, one must also match the $p^2$ term. This requires an additional
(local) operator in the $\EFT$ skin whose matrix element contributes a factor of $p^2$ relative to
$\mathcal{O}_6$. One convenient representative (basis-dependent, but sufficient for witness extraction) is
\begin{equation}\label{eq:O8}
\mathcal{O}_8 :=
\big(\bar\psi_1\gamma^\mu(1-\gamma_5)\psi_2\big)\,\Box\,
\big(\bar\psi_3\gamma_\mu(1-\gamma_5)\psi_4\big),
\end{equation}
up to integration by parts and equivalent operator-basis choices.
The corresponding Wilson coefficient is fixed by \eqref{eq:cc-expanded} as
\begin{equation}\label{eq:C8def}
C_8 = \frac{g_2^2}{8m_W^4}=\frac{C_6}{m_W^2}=\frac{G_F}{\sqrt2}\frac{1}{m_W^2}.
\end{equation}
Hence the ratio of the two invariant coefficients is
\begin{equation}\label{eq:ratio}
\frac{C_8}{C_6}=\frac{1}{m_W^2}.
\end{equation}
Within $S_{\CC}^{(8)}$, this ratio is an \emph{invariant witness}: it is extractable from the measured
momentum dependence of charged-current amplitudes at order $p^2/m_W^2$, and it is not determined by
$G_F$ alone. Therefore $m_W$ (equivalently $C_8$) is standing-relevant in the enlarged scope and cannot be
classified as redundancy there. This is precisely the mechanism described abstractly in
Remark~\ref{rem:scope-index}.
\end{proof}




\subsection{Electromagnetic collapse to $e$ (fully worked)}

\begin{lemma}[Electromagnetic collapse in the pure-photon scope]\label{lem:em}
Let $S_{\EM}^{(\gamma)}$ be the \emph{pure-photon electromagnetic scope} in which $Z$ exchange is
integrated out (or kinematically irrelevant) and electromagnetic observables are mediated by the
massless gauge field $A_\mu$ alone (Definition~\ref{def:scope}).
In $S_{\EM}^{(\gamma)}$, admissible transport from electroweak gauge bookkeeping to the QED skin
exposes a single invariant witness $e$. Consequently, treating $(g_1,g_2)$ as independent
\emph{electromagnetic} couplings overcounts tensor load for $S_{\EM}^{(\gamma)}$:
only the combination $e=\frac{g_1 g_2}{\sqrt{g_1^2+g_2^2}}$ has standing as electromagnetic strength in
that scope.
\end{lemma}

\begin{proof}
We again provide a transport certificate.

\paragraph{(C1) Input skin $\Skin_{\mathrm{EW}}$ (neutral gauge bookkeeping).}
For a fermion field $\psi$ with weak isospin generator $T^3$ and hypercharge $Y$ (convention
$Q=T^3+Y$),
\begin{equation}\label{eq:neutral-coupling}
\mathcal{L}\supset
\bar\psi\gamma^\mu\big(g_2 T^3 W^3_\mu + g_1 Y B_\mu\big)\psi.
\end{equation}

\paragraph{(C2) Neutral gauge-boson mass matrix from the Higgs kinetic term.}
Using the Higgs VEV as in Lemma~\ref{lem:cc}, the neutral part of $(D_\mu H)^\dagger(D^\mu H)$ gives
\begin{equation}\label{eq:neutral-mass}
(D_\mu H)^\dagger(D^\mu H)\supset \frac{v^2}{8}\,(g_2 W^3_\mu - g_1 B_\mu)^2.
\end{equation}
This quadratic form has one zero eigenvalue (massless photon) and one massive eigenstate ($Z$).

\paragraph{(C3) Diagonalization and definition of $\theta_W$.}
Define the Weinberg angle by
\begin{equation}\label{eq:thetaW}
\tan\theta_W=\frac{g_1}{g_2},\qquad
\sin\theta_W=\frac{g_1}{\sqrt{g_1^2+g_2^2}},\qquad
\cos\theta_W=\frac{g_2}{\sqrt{g_1^2+g_2^2}}.
\end{equation}
Then the orthogonal rotation
\begin{equation}\label{eq:rotAZ}
\begin{pmatrix} A_\mu \\ Z_\mu\end{pmatrix}
=
\begin{pmatrix}
\cos\theta_W & \sin\theta_W\\
-\sin\theta_W&\cos\theta_W
\end{pmatrix}
\begin{pmatrix} B_\mu\\ W^3_\mu\end{pmatrix}
\end{equation}
diagonalizes \eqref{eq:neutral-mass}, with $A_\mu$ massless.

\paragraph{(C4) Output skin $\Skin_{\mathrm{QED}}$ (electromagnetic bookkeeping).}
The QED skin records electromagnetic interactions as
\begin{equation}\label{eq:qed-skin}
\mathcal{L}_{\mathrm{QED}}\supset e\,A_\mu\,\bar\psi\gamma^\mu Q \psi,
\end{equation}
with a single coupling $e$ and charge operator $Q$.

\paragraph{(C5) Transport map and identification of the witness $e$.}
Insert \eqref{eq:rotAZ} into \eqref{eq:neutral-coupling}.
The coefficient of $A_\mu$ becomes
\begin{align}
g_2 T^3(\sin\theta_W) + g_1 Y(\cos\theta_W)
&=
(g_2\sin\theta_W)\,T^3+(g_1\cos\theta_W)\,Y.
\end{align}
Using \eqref{eq:thetaW}, define
\begin{equation}\label{eq:e-def}
e:=g_2\sin\theta_W=g_1\cos\theta_W=\frac{g_1 g_2}{\sqrt{g_1^2+g_2^2}}.
\end{equation}
Then the $A_\mu$ coupling is
\begin{equation}
e\,(T^3+Y)=eQ,
\end{equation}
which reproduces the QED form \eqref{eq:qed-skin}. Thus transport identifies \emph{one} electromagnetic
witness $e$.

\paragraph{(C6) Witness rank and collapse.}
The QED skin has a single coupling witness $e$ for electromagnetic interaction strength.
Electroweak bookkeeping uses $(g_1,g_2)$, but in the electromagnetic regime their only standing-relevant
combination is \eqref{eq:e-def}. Therefore $(g_1,g_2)$ overcounts tensor load for electromagnetic phenomena,
and Theorem~\ref{thm:witness-counting} forces collapse to $e$ as the unique electromagnetic witness.
\end{proof}

\begin{remark}[Scope warning for Lemma~\ref{lem:em}]\label{rem:em-scope}
Lemma~\ref{lem:em} is a statement about the pure-photon scope $S_{\EM}^{(\gamma)}$.
If the scope is enlarged to include $Z$ exchange, parity-violating neutral currents, or energies where
electroweak symmetry is effectively restored, then additional invariant witnesses (e.g.\ $\sin\theta_W$,
$m_Z$) become standing-relevant and the collapse to $e$ alone does not apply.
\end{remark}

\begin{proposition}[Worked example: enlarging $S_{\EM}^{(\gamma)}$ to include $Z$ exchange forces new witnesses]
\label{prop:em-Z}
Let $S_{\EM}^{(\gamma)}$ be the pure-photon scope of Lemma~\ref{lem:em}.
Define an enlarged scope $S_{\EM}^{(\gamma Z,6)}$ that additionally admits the leading ($q^2\ll m_Z^2$)
neutral-current correction from $Z$ exchange in scattering of charged fermions, including polarization-
dependent observables. Then the witness set is strictly enlarged: beyond $e$, a new invariant witness is
forced, and one may take it to be $s_W^2\equiv\sin^2\theta_W$.
If the scope is further enlarged to $S_{\EM}^{(\gamma Z,8)}$ which resolves the first propagator correction
at order $q^2/m_Z^2$, then an additional invariant witness is forced, and one may take it to be $m_Z$.
\end{proposition}

\begin{proof}
We exhibit the mechanism explicitly, in the same certificate style as Proposition~\ref{prop:cc-p2}.

\paragraph{(E1) $Z$ coupling and current structure.}
From the rotation \eqref{eq:rotAZ}, the $Z$ coupling to a fermion with generators $(T^3,Q)$ is
\begin{equation}\label{eq:Z-coupling}
\mathcal{L}\supset g_Z\,Z_\mu\,\bar\psi\gamma^\mu\big(T^3-s_W^2 Q\big)\psi,
\qquad g_Z:=\frac{g_2}{\cos\theta_W}=\sqrt{g_1^2+g_2^2},\qquad s_W\equiv\sin\theta_W.
\end{equation}
In particular, for the electron ($Q_e=-1$) one has
\begin{equation}\label{eq:e-chiral-Z}
g_L^e = T^3_{e_L}-s_W^2 Q_e=-\frac12+s_W^2,\qquad
g_R^e = T^3_{e_R}-s_W^2 Q_e=0+s_W^2=s_W^2.
\end{equation}

\paragraph{(E2) $Z$-exchange amplitude and low-energy expansion.}
Consider $t$-channel exchange in (for definiteness) electron--electron scattering with momentum transfer
$q^2$ (spacelike in elastic scattering). The $Z$-exchange contribution has the schematic form
\begin{equation}\label{eq:MZschem}
\mathcal{M}_Z(q^2)\;\propto\;-
\frac{g_Z^2}{m_Z^2-q^2}\,
\big[g_L^e\,\big(\bar e_L\gamma^\mu e_L\big)+g_R^e\,\big(\bar e_R\gamma^\mu e_R\big)\big]^2,
\end{equation}
where $m_Z$ is the $Z$ mass. For $|q^2|\ll m_Z^2$, expand
\begin{equation}\label{eq:Zprop-expand}
\frac{1}{m_Z^2-q^2}=\frac{1}{m_Z^2}\left(1+\frac{q^2}{m_Z^2}+\cO\!\left(\frac{q^4}{m_Z^4}\right)\right).
\end{equation}

\paragraph{(E3) Matching at dimension six: $s_W^2$ becomes a witness.}
Keeping only the leading term of \eqref{eq:Zprop-expand} (the scope $S_{\EM}^{(\gamma Z,6)}$), the $Z$ exchange
is represented in an EFT skin by local four-fermion operators. A convenient (chiral) basis is
\begin{equation}\label{eq:O6em}
\mathcal{O}^{ee}_{6,ab}:=\big(\bar e_a\gamma^\mu e_a\big)\big(\bar e_b\gamma_\mu e_b\big),\qquad a,b\in\{L,R\},
\end{equation}
with Wilson coefficients read off from \eqref{eq:MZschem} at order $q^0/m_Z^2$:
\begin{equation}\label{eq:C6em}
C^{ee}_{6,ab}=\frac{g_Z^2}{m_Z^2}\,g_a^e g_b^e\qquad (a,b\in\{L,R\}).
\end{equation}
Crucially, the \emph{relative} chiral strengths depend on $s_W^2$ through \eqref{eq:e-chiral-Z}. For example,
\begin{equation}\label{eq:ratio-sW}
\frac{C^{ee}_{6,RR}}{C^{ee}_{6,LL}}=\left(\frac{g_R^e}{g_L^e}\right)^2
=\left(\frac{s_W^2}{s_W^2-\tfrac12}\right)^2,
\end{equation}
which is independent of the overall prefactor $g_Z^2/m_Z^2$.
Within $S_{\EM}^{(\gamma Z,6)}$ (where polarization-dependent observables are admitted), the left--left and
right--right neutral-current contact strengths are separately measurable in principle, so the ratio
\eqref{eq:ratio-sW} is an \emph{invariant witness}. It is \emph{not} determined by the pure-photon witness $e$,
since $e$ does not fix $s_W^2$ (Lemma~\ref{lem:em} only identifies the $A_\mu$ coupling).
Thus enlarging $S_{\EM}^{(\gamma)}$ to $S_{\EM}^{(\gamma Z,6)}$ forces at least one new witness; a convenient
choice is $s_W^2$.

\paragraph{(E4) Matching at dimension eight: $m_Z$ becomes a witness.}
If the scope is enlarged further to $S_{\EM}^{(\gamma Z,8)}$ so that the $q^2/m_Z^2$ term in
\eqref{eq:Zprop-expand} is within the admitted accuracy, then additional (derivative) operators are required.
A representative choice mirroring \eqref{eq:O8} is
\begin{equation}\label{eq:O8em}
\mathcal{O}^{ee}_{8,ab}:=\big(\bar e_a\gamma^\mu e_a\big)\,\Box\,\big(\bar e_b\gamma_\mu e_b\big)
\qquad (a,b\in\{L,R\}).
\end{equation}
Matching the $q^2$ term yields
\begin{equation}\label{eq:C8em}
C^{ee}_{8,ab}=\frac{g_Z^2}{m_Z^4}\,g_a^e g_b^e.
\end{equation}
Therefore, for any fixed chiral channel $(a,b)$,
\begin{equation}\label{eq:ratio-mZ}
\frac{C^{ee}_{8,ab}}{C^{ee}_{6,ab}}=\frac{1}{m_Z^2}.
\end{equation}
In $S_{\EM}^{(\gamma Z,8)}$, the momentum dependence of the neutral-current correction determines the ratio
\eqref{eq:ratio-mZ}, hence $m_Z$ is a standing-relevant invariant witness there.

\paragraph{Conclusion.}
The $Z$ contribution is not ``redundancy suppressed by heaviness'': it is simply outside
$S_{\EM}^{(\gamma)}$. Once the scope is enlarged to admit the corresponding observables (contact neutral currents
and/or their $q^2$ dependence), additional witnesses are forced by explicit matching, exactly as in
Proposition~\ref{prop:cc-p2}.
\end{proof}

\begin{proposition}[Worked example: enlarging a pure-QED photon scope to include hadronic vacuum polarization forces new witnesses]
\label{prop:em-hvp}
Let $S_{\EM}^{(\gamma,\ell)}$ denote a \emph{pure-photon QED scope} in which electromagnetic processes are mediated by the photon alone and vacuum-polarization effects are admitted only from the \emph{leptonic} sector (so that the photon vacuum polarization is perturbatively computable within QED, given $e$ and lepton masses).
Define an enlarged scope $S_{\EM}^{(\gamma,\ell+\mathrm{had})}$ that additionally admits hadronic vacuum-polarization (HVP) corrections to photon exchange at the same perturbative order.
Then the witness set is strictly enlarged: beyond $e$, additional invariant witnesses are forced.
If the enlarged scope admits only a \emph{finite} family of HVP-sensitive observables, a finite witness set may be taken as the corresponding finite list of spectral integrals.
If the enlarged scope admits energy-resolved hadronic production (so that $R_{\mathrm{had}}(s)$ itself is an observable function of $s$) or otherwise a separating family of kernels, the witness is naturally represented by the hadronic spectral function $R_{\mathrm{had}}(s)$ (equivalently the subtracted hadronic vacuum polarization $\Pi_{\mathrm{had}}(q^2)-\Pi_{\mathrm{had}}(0)$), which enters electromagnetic observables through explicit dispersion integrals.
\end{proposition}

\begin{proof}
We exhibit the scope enlargement explicitly and identify the new witnesses.

\paragraph{(H1) Photon vacuum polarization and the running coupling as a witness.}
Let $j_{\rm em}^\mu$ be the electromagnetic current (defined so that the interaction is $\mathcal{L}_{\rm int}=e A_\mu j_{\rm em}^\mu$).
Define the (transverse) current correlator
\begin{equation}\label{eq:Pi-def}
\Pi^{\mu\nu}(q):= i\int d^4x\,e^{iq\cdot x}\,\langle 0|T\,j_{\rm em}^\mu(x)j_{\rm em}^\nu(0)|0\rangle
=\big(q^\mu q^\nu-q^2 g^{\mu\nu}\big)\,\Pi(q^2).
\end{equation}
The photon self-energy is $e^2\Pi^{\mu\nu}$, so resumming vacuum-polarization insertions yields (in a convenient gauge)
\begin{equation}\label{eq:photon-resummed}
D_{\mu\nu}(q)=\frac{-ig_{\mu\nu}}{q^2\big(1-\widehat\Pi(q^2)\big)}+\cdots,
\qquad \widehat\Pi(q^2):=-e^2\big(\Pi(q^2)-\Pi(0)\big),
\end{equation}
where the subtraction at $q^2=0$ fixes the on-shell normalization and ``$\cdots$'' are gauge terms that do not affect conserved currents.
Accordingly, photon-mediated amplitudes acquire the replacement
\begin{equation}\label{eq:alpha-running}
\frac{e^2}{q^2}\quad\longmapsto\quad\frac{e^2}{q^2\big(1-\widehat\Pi(q^2)\big)}\equiv \frac{e^2_{\rm eff}(q^2)}{q^2},
\qquad e^2_{\rm eff}(q^2):=\frac{e^2}{1-\widehat\Pi(q^2)}.
\end{equation}
Thus, once vacuum polarization is admitted in the scope, the \emph{function} $\widehat\Pi(q^2)$ (equivalently the running coupling $e_{\rm eff}(q^2)$) is an invariant witness: it is extractable from momentum-dependent electromagnetic observables and cannot be removed by admissible redescription.

\paragraph{(H2) Leptonic-only scope $S_{\EM}^{(\gamma,\ell)}$.}
In the leptonic-only scope, $\widehat\Pi(q^2)=\widehat\Pi_{\ell}(q^2)$ is perturbatively computable in QED in terms of $e$ and lepton masses.
No additional, nonperturbative witness is required.

\paragraph{(H3) Enlarged scope $S_{\EM}^{(\gamma,\ell+\mathrm{had})}$ and extraction of a new witness.}
In the enlarged scope, the vacuum polarization splits as
\begin{equation}\label{eq:Pi-split}
\widehat\Pi(q^2)=\widehat\Pi_{\ell}(q^2)+\widehat\Pi_{\mathrm{had}}(q^2).
\end{equation}
Consider, for definiteness, $e^+e^-\to\mu^+\mu^-$ at center-of-mass energy $s=q^2$ in a regime where $Z$ exchange is outside the scope (so the amplitude is photon-dominated).
At leading order in the resummed vacuum polarization, the total cross section can be written as
\begin{equation}\label{eq:sigma-mumu-running}
\sigma_{\mu\mu}(s)=\sigma^{(0)}_{\mu\mu}(s)\,\left|\frac{1}{1-\widehat\Pi(s)}\right|^2,
\qquad \sigma^{(0)}_{\mu\mu}(s)=\frac{4\pi\alpha^2}{3s},\qquad \alpha:=\frac{e^2}{4\pi}.
\end{equation}
Therefore the measured ratio
\begin{equation}\label{eq:Pi-extract}
\left(\frac{\sigma_{\mu\mu}(s)}{\sigma^{(0)}_{\mu\mu}(s)}\right)^{1/2}-1
\;\approx\;\Re\,\widehat\Pi(s)
\end{equation}
extracts $\widehat\Pi(s)$ as an invariant witness (to the admitted perturbative order).
Since $\widehat\Pi_{\ell}(s)$ is computable in QED, the hadronic contribution
\begin{equation}\label{eq:Pi-had-extract}
\widehat\Pi_{\mathrm{had}}(s)=\widehat\Pi(s)-\widehat\Pi_{\ell}(s)
\end{equation}
is itself an invariant witness in the enlarged scope.
Crucially, $\widehat\Pi_{\mathrm{had}}$ is \emph{not} determined by $e$ (or by the leptonic data) within the pure-QED skin.

\paragraph{(H4) Spectral representation: the hadronic witness is a dispersion integral.}
Analyticity of the correlator \eqref{eq:Pi-def} implies a subtracted dispersion relation for the scalar function $\Pi(q^2)$:
\begin{equation}\label{eq:dispersion}
\Pi(q^2)-\Pi(0)=\frac{q^2}{\pi}\int_{s_{\rm th}}^{\infty}ds\,\frac{\Im\Pi(s)}{s\,(s-q^2-i\epsilon)},
\end{equation}
where $s_{\rm th}$ is the hadronic threshold.
By the optical theorem, the hadronic spectral function is determined by the ratio
\begin{equation}\label{eq:Rdef}
R_{\mathrm{had}}(s):=\frac{\sigma(e^+e^-\to\mathrm{hadrons};s)}{\sigma^{(0)}_{\mu\mu}(s)},\qquad \sigma^{(0)}_{\mu\mu}(s)=\frac{4\pi\alpha^2}{3s},
\end{equation}
and one has the standard relation
\begin{equation}\label{eq:ImPi-R}
R_{\mathrm{had}}(s)=12\pi\,\Im\Pi_{\mathrm{had}}(s).
\end{equation}
Inserting \eqref{eq:ImPi-R} into \eqref{eq:dispersion} and multiplying by $-e^2$ yields the explicit spectral integral for the hadronic witness:
\begin{equation}\label{eq:DeltaAlphaHVP}
\widehat\Pi_{\mathrm{had}}(q^2)
=-e^2\,\Re\big(\Pi_{\mathrm{had}}(q^2)-\Pi_{\mathrm{had}}(0)\big)
=-\frac{\alpha q^2}{3\pi}\,\mathrm{PV}\!\int_{s_{\rm th}}^{\infty}ds\,\frac{R_{\mathrm{had}}(s)}{s\,(s-q^2)}.
\end{equation}
Equation \eqref{eq:DeltaAlphaHVP} exhibits the new witness explicitly as a \emph{spectral integral}.

\paragraph{(H5) Kernel families and witness dimension.}
More generally, in any enlarged electromagnetic scope $S$ admitting HVP-sensitive observables whose hadronic contribution depends linearly on $R_{\mathrm{had}}$, one may write
\begin{equation}\label{eq:kernel-functional}
\Delta \mathcal{O}_K \;=\;\int_{s_{\rm th}}^{\infty} ds\,K(s)\,R_{\mathrm{had}}(s),\qquad K\in\mathcal{K}_S,
\end{equation}
for a set of admitted kernels $\mathcal{K}_S$ determined by which observables are included in $S$.
The kernel set $\mathcal{K}_S$ fixes the minimal witness requirement:
\begin{enumerate}[label=(\roman*),leftmargin=2.2em]
\item If $\mathcal{K}_S$ is finite, then a witness set may be taken as the finite list of numbers
$\big\{\int ds\,K_i(s)\,R_{\mathrm{had}}(s)\big\}$ associated with those kernels.
\item If $\mathcal{K}_S$ \emph{separates} spectral functions on the admissible class (i.e.\ whenever
$R\neq R'$ there exists $K\in\mathcal{K}_S$ with $\int K R\neq \int K R'$), then $R_{\mathrm{had}}(s)$ itself is forced as a standing-relevant witness (and the witness rank is infinite in the sense of Remark~\ref{rem:witness-rank-functional}).
In particular, if the scope admits the energy-resolved ratio $R_{\mathrm{had}}(s)$ of \eqref{eq:Rdef} as an observable function of $s$ (as is typical once HVP is admitted beyond a single integrated observable), then the separating condition holds trivially and a functional witness is unavoidable.
\end{enumerate}
Thus, admitting HVP effects enlarges witness rank: at minimum it adds the required spectral integrals, and in the energy-resolved scope it forces the full function $R_{\mathrm{had}}(s)$ (equivalently $\Pi_{\mathrm{had}}(q^2)$), not merely the single coupling witness $e$.

\paragraph{Conclusion.}
Hadronic vacuum polarization is not ``redundancy suppressed by nonperturbativity''; it is outside the leptonic-only scope.
Once the scope is enlarged so that HVP-sensitive observables are within the admitted accuracy, new invariant witnesses are forced and are naturally represented by spectral integrals such as \eqref{eq:DeltaAlphaHVP}.
\end{proof}

\begin{remark}[Scope interface: QED bookkeeping vs.\ QCD content]\label{rem:hvp-scope}
Proposition~\ref{prop:em-hvp} is a scope statement about what a \emph{pure-QED electromagnetic skin} must import once HVP effects are admitted.
Within the full Standard Model declared load, the hadronic spectral function $R_{\mathrm{had}}(s)$ is not posited as a new primitive; it is a witness of the QCD sector and can, in principle, be obtained either by direct measurement (via \eqref{eq:Rdef}) or by QCD computation.
The point here is that it is \emph{not determined by $e$ alone} within the pure-photon QED skin, so enlarging the electromagnetic scope forces new witnesses and therefore cannot be misclassified as ``redundancy.''
\end{remark}





\subsection{Yukawa--mass equivalence (fully worked, scope-explicit)}

\begin{lemma}[Yukawa--mass equivalence, per species]\label{lem:yukawa-mass}
Fix a fermion species $f$ and work within the declared load.
In the broken-phase skin, the Yukawa coupling $y_f$ and the electroweak scale $v$ do not represent
independent standing-relevant degrees: they enter observables only through
\begin{equation}\label{eq:mf}
m_f=\frac{y_f v}{\sqrt2},
\end{equation}
and (when Higgs interactions are kept) through the ratio $m_f/v$.
Consequently, $y_f$ has no independent standing beyond $(m_f,v)$, and in the low-energy \emph{mass-only}
scope the unique witness is $m_f$.
\end{lemma}

\begin{proof}
Write the Yukawa operator (schematically; $H$ or $\tilde H$ as appropriate)
\begin{equation}\label{eq:yukawa}
\mathcal{L}_Y=-y_f\,\bar f_L H f_R+\mathrm{h.c.}
\end{equation}
Insert $H=\frac{1}{\sqrt2}(0,v+h)^T$ to obtain
\begin{equation}\label{eq:yukawa-expanded}
\mathcal{L}_Y=-\frac{y_f v}{\sqrt2}\,\bar f f -\frac{y_f}{\sqrt2}\,h\,\bar f f.
\end{equation}
The first term is the fermion mass term, hence \eqref{eq:mf}.
The second term is the Higgs--fermion coupling; using \eqref{eq:mf} it becomes
\begin{equation}
-\frac{m_f}{v}\,h\,\bar f f.
\end{equation}
Thus, in any broken-phase skin that retains the Higgs field $h$, the Yukawa data appear only as
$(m_f,v)$ rather than $(y_f,v)$.
In any low-energy scope where Higgs exchange is integrated out (or irrelevant), standing-relevant
fermion-sector content is witnessed by $m_f$ alone (thresholds, propagators, kinematics), so $y_f$
has no independent standing for that restricted scope.

Formally: the input parameterization $(y_f,v)$ has rank $2$, while the witness set for the mass-only
scope has rank $1$ ($m_f$). By Theorem~\ref{thm:witness-counting}, $(y_f,v)$ cannot be two independent
standing-relevant degrees in that scope.
\end{proof}

\begin{remark}[Interpretation warning]\label{rem:yuk-interpret}
Lemma~\ref{lem:yukawa-mass} does not deny that Higgs--fermion couplings are measurable. It states that,
within the declared load, such measurements witness the \emph{quotient combinations} $m_f$ and $m_f/v$,
not an additional standing degree beyond those witnesses for the chosen scope.
\end{remark}


\subsection{QCD CP--odd collapse to $\bar\theta$ (fully worked)}

\begin{lemma}[Strong CP: unique invariant $\bar\theta$]\label{lem:theta}
In the QCD CP--odd sector, admissible chiral rephasing collapses apparent independence of the
$\theta$ term and quark-mass phases. The unique standing-relevant invariant is
\begin{equation}\label{eq:thetabar}
\bar\theta:=\theta+\arg\det M,
\end{equation}
where $M$ is the quark mass matrix.
\end{lemma}

\begin{proof}
We show explicitly how admissible rephasing acts, and compute the invariant.

\paragraph{(1) CP--odd operators in the ultraviolet skin.}
The relevant terms are
\begin{equation}\label{eq:qcd-cpodd}
\mathcal{L}\supset
-\bar q_R M q_L -\bar q_L M^\dagger q_R
+ \frac{g_s^2}{32\pi^2}\,\theta\,G^a_{\mu\nu}\tilde G^{a\mu\nu}.
\end{equation}

\paragraph{(2) Global axial rephasing of all quark flavors.}
Consider the (classically admissible) field redefinition
\begin{equation}\label{eq:axial}
q\mapsto e^{i\alpha\gamma_5}q,
\qquad
q_L\mapsto e^{-i\alpha}q_L,\quad q_R\mapsto e^{+i\alpha}q_R.
\end{equation}
Under \eqref{eq:axial}, the mass term transforms as
\begin{equation}
-\bar q_R M q_L \;\mapsto\; -\bar q_R e^{-2i\alpha} M q_L,
\end{equation}
so
\begin{equation}\label{eq:argdet-shift}
\arg\det M\;\mapsto\;\arg\det M -2N_f\alpha,
\end{equation}
where $N_f$ is the number of quark flavors rotated.

\paragraph{(3) Anomalous Jacobian and the $\theta$ shift.}
In the quantum theory, the transformation \eqref{eq:axial} changes the path-integral measure
(Fujikawa) and induces an additional term in the effective action:
\begin{equation}\label{eq:theta-shift}
\theta\;\mapsto\;\theta+2N_f\alpha.
\end{equation}
(Equivalently: the axial current has an anomaly proportional to $G\tilde G$, and the chiral rotation
shifts the coefficient of $G\tilde G$.)

\paragraph{(4) Invariant combination.}
Combining \eqref{eq:argdet-shift} and \eqref{eq:theta-shift} gives
\begin{equation}
\theta+\arg\det M\;\mapsto\;(\theta+2N_f\alpha)+(\arg\det M-2N_f\alpha)=\theta+\arg\det M.
\end{equation}
Thus $\bar\theta$ in \eqref{eq:thetabar} is invariant under admissible chiral rephasing.
Neither $\theta$ nor $\arg\det M$ is separately invariant; hence they do not have independent standing.

\paragraph{(5) Witness rank.}
Infrared CP--odd hadronic observables depend on a single parameter in this sector; $\bar\theta$ is the
unique candidate invariant under admissible rephasing.
Therefore this scope has rank $1$ witness content, forcing collapse to $\bar\theta$ as claimed.
\end{proof}

\subsection{Neutrino oscillation invariants (fully worked)}

\begin{lemma}[Oscillation witness set]\label{lem:osc}
In the neutrino oscillation regime, transition probabilities depend only on:
\begin{enumerate}[label=(\roman*),leftmargin=2.2em]
\item mass-squared differences $\Delta m^2_{ij}=m_i^2-m_j^2$,
\item rephasing-invariant products $U_{\alpha i}U_{\beta i}^\ast U_{\alpha j}^\ast U_{\beta j}$.
\end{enumerate}
Absolute mass offsets and rephasing-removable phases are redundancy for oscillation phenomena.
\end{lemma}

\begin{proof}
Let $U$ be the mixing matrix relating flavor and mass eigenstates:
\begin{equation}\label{eq:mix}
\ket{\nu_\alpha}=\sum_i U_{\alpha i}^\ast \ket{\nu_i}.
\end{equation}
In the relativistic limit, propagation of $\ket{\nu_i}$ over baseline $L$ with energy $E$ yields a phase
$\exp(-i m_i^2 L/(2E))$ (up to an overall common phase that cancels in probabilities). Therefore the
transition amplitude is
\begin{equation}\label{eq:amp}
\mathcal{A}_{\alpha\to\beta}
=
\sum_i U_{\beta i}\,e^{-i m_i^2 L/(2E)}\,U_{\alpha i}^\ast.
\end{equation}
The probability is
\begin{align}\label{eq:probfull}
P_{\alpha\to\beta}
&=
|\mathcal{A}_{\alpha\to\beta}|^2
=
\sum_{i,j}U_{\beta i}U_{\alpha i}^\ast U_{\beta j}^\ast U_{\alpha j}\,
\exp\!\left(-i\frac{(m_i^2-m_j^2)L}{2E}\right)\\
&=
\sum_{i,j}U_{\beta i}U_{\alpha i}^\ast U_{\beta j}^\ast U_{\alpha j}\,
\exp\!\left(-i\frac{\Delta m^2_{ij}L}{2E}\right).\nonumber
\end{align}
This exhibits the witness set directly: only $\Delta m^2_{ij}$ appear, and only the quartic products
of $U$ appear.

Now check admissible rephasing invariance. Rephase flavor and mass states by
\begin{equation}
\ket{\nu_\alpha}\mapsto e^{i\varphi_\alpha}\ket{\nu_\alpha},
\qquad
\ket{\nu_i}\mapsto e^{i\psi_i}\ket{\nu_i}.
\end{equation}
Then $U_{\alpha i}\mapsto e^{i\varphi_\alpha}U_{\alpha i}e^{-i\psi_i}$, and the quartic product in
\eqref{eq:probfull} is invariant:
\begin{equation}
U_{\beta i}U_{\alpha i}^\ast U_{\beta j}^\ast U_{\alpha j}
\mapsto
(e^{i\varphi_\beta}U_{\beta i}e^{-i\psi_i})
(e^{-i\varphi_\alpha}U_{\alpha i}^\ast e^{i\psi_i})
(e^{-i\varphi_\beta}U_{\beta j}^\ast e^{i\psi_j})
(e^{i\varphi_\alpha}U_{\alpha j}e^{-i\psi_j})
=
U_{\beta i}U_{\alpha i}^\ast U_{\beta j}^\ast U_{\alpha j}.
\end{equation}
Finally, shifting $m_i^2\mapsto m_i^2+c$ leaves all $\Delta m^2_{ij}$ unchanged and thus produces no
new witness.

Therefore the oscillation regime witnesses only $\Delta m^2_{ij}$ and the rephasing-invariant quartic
products, as claimed.
\end{proof}

\subsection{Longitudinal--Goldstone equivalence (worked proof skeleton)}

\begin{lemma}[Longitudinal amplitude-class collapse]\label{lem:LG}
In the electroweak high-energy regime $E\gg m_W$, on-shell transport identifies longitudinal vector
boson amplitudes with Goldstone amplitudes up to $\cO(m_W/E)$ corrections. Consequently, treating
longitudinal polarizations and Goldstone modes as independent tensor degrees overcounts
standing-relevant longitudinal content in that regime.
\end{lemma}

\begin{proof}
We provide the minimal verifiable chain (the full BRST proof is standard; the steps below isolate the
standing-relevant content).

\paragraph{(1) On-shell skin and longitudinal polarization.}
Consider an amplitude with an external massive gauge boson of momentum $p^\mu$.
A longitudinal polarization vector can be chosen to satisfy
\begin{equation}\label{eq:epsL}
\epsilon_L^\mu(p)=\frac{p^\mu}{m_W}+\cO\!\left(\frac{m_W}{E}\right),
\qquad E\sim |p^0|\gg m_W.
\end{equation}

\paragraph{(2) Ward/Slavnov--Taylor identity in a broken gauge theory.}
In $R_\xi$ gauges, BRST invariance yields a Slavnov--Taylor identity relating the contraction of the
(amputated) gauge amplitude with $p_\mu$ to the amplitude with the corresponding Goldstone field $\pi$.
A derivation in the present conventions is given in Appendix~\ref{app:LG-proof}.
For the standing-relevant content we need, it suffices to record the identity in the form
\begin{equation}\label{eq:ward}
p_\mu\,\mathcal{M}^\mu(W^a(p),\dots)= i\,m_W\,\mathcal{M}(\pi^a(p),\dots),
\end{equation}
up to terms that vanish for physical on-shell external states. The overall phase is a convention
artifact of LSZ normalization and plays no role in the amplitude-class identification used here.
\paragraph{(3) High-energy reduction.}
Contract the gauge amplitude with the longitudinal polarization:
\begin{equation}
\mathcal{M}(W_L^a(p),\dots)
=
\epsilon_{L\mu}(p)\,\mathcal{M}^\mu(W^a(p),\dots)
=
\left(\frac{p_\mu}{m_W}+\cO\!\left(\frac{m_W}{E}\right)\right)\mathcal{M}^\mu(W^a(p),\dots).
\end{equation}
Using \eqref{eq:ward} gives
\begin{equation}
\mathcal{M}(W_L^a(p),\dots)
=
i\,\mathcal{M}(\pi^a(p),\dots)+\cO\!\left(\frac{m_W}{E}\right).
\end{equation}

\paragraph{(4) Standing consequence.}
In the high-energy on-shell regime, the amplitude class is witnessed by $\mathcal{M}$; the above shows
that the longitudinal gauge-boson amplitude and Goldstone amplitude are the \emph{same witness} up to a convention-dependent overall phase (here $i$) and
controlled regime-suppressed corrections. Therefore counting $W_L$ and $\pi$ as independent tensor
degrees overcounts standing-relevant content for that regime.
\end{proof}

\subsection{Parameter normal form (global statement)}

\begin{proposition}[Parameter normal form: witnesses exhaust standing]\label{prop:normal-form}
Fix the declared load and a scope. A parameter has standing as an independent tensor degree for that
scope only if it corresponds to an independent invariant witness across all admissible skins and
survives admissible transport. Any parameter not so witnessed is redundancy for that scope, unless
its retention can be justified without invoking standing-failure modes.
\end{proposition}

\section{Forced discrete skeleton: operators, anomalies, charges, FCNC}

\subsection{Operator existence is transport-invariant}

\begin{lemma}[Operator existence invariance]\label{lem:operator}
Within the declared load, existence of a broken-phase mass pairing for a chiral fermion is equivalent
to existence of a corresponding gauge-invariant Yukawa operator in the unbroken-phase skin.
Absent the operator, no admissible redescription can create the pairing without forbidden patches.
\end{lemma}

\begin{proof}
In the declared load, the renormalizable Lagrangian in the unbroken phase contains only gauge-invariant
operators built from the available fields.
Electroweak symmetry breaking is a \emph{substitution} (choose a gauge and expand $H$ around its VEV),
followed by admissible field redefinitions (diagonalization, basis change). Such operations cannot
create a term that is not the image of an existing gauge-invariant operator.

Therefore if a broken-phase mass term $\bar f f$ appears, it must arise from some gauge-invariant
operator $\bar f_L H f_R$ or $\bar f_L\tilde H f_R$ already present (possibly after field relabeling).
Conversely, if the corresponding gauge-invariant Yukawa operator is absent, producing the mass pairing
requires adding (i) a new Higgs multiplet with different charges, (ii) new fermions enabling an
effective pairing, or (iii) non-gauge-invariant/higher-dimensional patch operators. Each is forbidden
by the declared load (Definition~\ref{def:declared-load}). Hence operator existence is transport-invariant.
\end{proof}

\subsection{Anomaly cancellation and Yukawa compatibility (full algebra)}\label{sec:anom-yuk}

We now show, explicitly, how the hypercharge structure is forced under the declared load.
This subsection uses a \emph{single consistent convention} throughout:
\[
Q = T^3 + Y,
\]
and the Higgs VEV is taken in the electrically neutral component of the doublet
$H=(0,(v+h)/\sqrt2)^T$, which fixes $Y_H$ relative to $T^3$ by $Q(H^0)=0$.

\subsubsection{One-generation notation (right-handed hypercharges)}
Work with one generation of fermions in the standard representation skeleton, written in terms of
their usual chiralities:
\begin{center}
\begin{tabular}{@{}lll@{}}
\toprule
Field & $\SU(3)_c\times \SU(2)_L$ rep & Hypercharge \\ \midrule
$Q_L$ (quark doublet) & $(\mathbf{3},\mathbf{2})$ & $Y_Q$ \\
$u_R$ (up singlet) & $(\mathbf{3},\mathbf{1})$ & $Y_u$ \\
$d_R$ (down singlet) & $(\mathbf{3},\mathbf{1})$ & $Y_d$ \\
$L_L$ (lepton doublet) & $(\mathbf{1},\mathbf{2})$ & $Y_L$ \\
$e_R$ (charged-lepton singlet) & $(\mathbf{1},\mathbf{1})$ & $Y_e$ \\
\bottomrule
\end{tabular}
\end{center}
The Higgs doublet has hypercharge $Y_H$ and its conjugate $\tilde H=i\sigma^2 H^\ast$ has hypercharge
$-Y_H$.

\subsubsection{Yukawa gauge-invariance conditions}
The renormalizable Yukawa terms in the declared load are
\begin{equation}\label{eq:yukawas}
\mathcal{L}_Y\supset
-\bar Q_L\,Y_u\,\tilde H\,u_R
-\bar Q_L\,Y_d\,H\,d_R
-\bar L_L\,Y_e\,H\,e_R
+\mathrm{h.c.}
\end{equation}
Gauge invariance under $\U(1)_Y$ fixes the hypercharge relations
\begin{equation}\label{eq:yuk-rel}
Y_u = Y_Q + Y_H,\qquad
Y_d = Y_Q - Y_H,\qquad
Y_e = Y_L - Y_H.
\end{equation}
(These are the familiar Standard-Model relations in the convention $Q=T^3+Y$.)

\subsubsection{Anomaly cancellation conditions (one generation)}
Local anomaly constraints in chiral gauge theories are standard; see e.g.\ \cite{Bertlmann1996,WeinbergQTF2}.
Anomaly cancellation for a gauged $\U(1)_Y$ requires vanishing of:
\begin{align}
[\SU(3)]^2\U(1):\quad & 2Y_Q - Y_u - Y_d =0, \label{eq:A3}\\
[\SU(2)]^2\U(1):\quad & 3Y_Q + Y_L=0, \label{eq:A2}\\
[\text{grav}]^2\U(1):\quad & 6Y_Q + 2Y_L - 3Y_u - 3Y_d - Y_e =0, \label{eq:Agrav}\\
[\U(1)]^3:\quad & 6Y_Q^3+2Y_L^3-3Y_u^3-3Y_d^3-Y_e^3=0. \label{eq:A1}
\end{align}
The minus signs for right-handed fields appear because anomalies are computed for left-handed Weyl
fermions, so $u_R,d_R,e_R$ contribute with opposite hypercharge.

\begin{lemma}[Global $\SU(2)$ (Witten) anomaly check]\label{lem:witten}
In addition to the local triangle anomalies \eqref{eq:A3}--\eqref{eq:A1}, a consistent $\SU(2)_L$ gauge theory of
left-handed Weyl fermions must have an even number of $\SU(2)$ doublets (the Witten $\mathbb{Z}_2$ anomaly)
\cite{Witten1982}. In the minimal chiral skeleton, one generation contains $3$ quark doublets (one per color) plus
one lepton doublet, for a total of $4$ doublets; hence the global $\SU(2)$ anomaly cancels (and for $n_g$ generations,
$4n_g$ remains even).
\end{lemma}

\begin{proof}
The nontrivial element of $\pi_4(\SU(2))\cong\mathbb{Z}_2$ can flip the sign of the fermion determinant for an odd
number of left-handed doublets; consistency requires the sign be trivial. Counting $\SU(2)$ doublets in the minimal
skeleton gives $3+1=4$ per generation, hence even.
\end{proof}



\begin{proposition}[Constraint coupling inside the declared skeleton]\label{prop:constraint-coupling}
Within the minimal chiral skeleton:
\begin{enumerate}[label=(\roman*),leftmargin=2.2em]
\item If the Yukawa relations \eqref{eq:yuk-rel} hold, then the $[\SU(3)]^2\U(1)$ anomaly
\eqref{eq:A3} cancels identically.
\item If the Yukawa relations \eqref{eq:yuk-rel} hold and the linear anomalies \eqref{eq:A2} and
\eqref{eq:Agrav} vanish, then the cubic anomaly \eqref{eq:A1} vanishes automatically.
\end{enumerate}
\end{proposition}

\begin{proof}
(i) Substitute \eqref{eq:yuk-rel} into \eqref{eq:A3}:
\[
2Y_Q-(Y_Q+Y_H)-(Y_Q-Y_H)=0.
\]

(ii) Under \eqref{eq:yuk-rel}, the mixed gravitational anomaly becomes
\[
6Y_Q+2Y_L-3(Y_Q+Y_H)-3(Y_Q-Y_H)-(Y_L-Y_H)=Y_L+Y_H.
\]
Thus \eqref{eq:Agrav}=0 forces $Y_L=-Y_H$.
Together with \eqref{eq:A2} (which gives $Y_L=-3Y_Q$) we obtain $Y_H=3Y_Q$.
Substituting $Y_L=-3Y_Q$ and $Y_H=3Y_Q$ into \eqref{eq:yuk-rel} gives
$Y_u=4Y_Q$, $Y_d=-2Y_Q$, $Y_e=-6Y_Q$.
A direct substitution of these relations into \eqref{eq:A1} gives $0$ (see
Appendix~\ref{app:cubic-check} for the explicit algebra).
\end{proof}

\begin{theorem}[Discrete-skeleton hypercharge collapse]\label{thm:anomaly-yukawa}
Within the declared load (minimal chiral skeleton + one Higgs doublet with the charged Yukawa operator content
\eqref{eq:yukawas}), anomaly cancellation forces the one-generation hypercharges to the one-parameter family
\begin{equation}\label{eq:hyper-solution}
Y_L=-3Y_Q,\qquad
Y_e=-6Y_Q,\qquad
Y_H=3Y_Q,\qquad
Y_u=4Y_Q,\qquad
Y_d=-2Y_Q,
\end{equation}
unique up to an overall normalization (absorbed into $g_1$) and an overall sign (charge conjugation).

Equivalently: the local anomaly equations alone determine $Y_L=-3Y_Q$, $Y_e=-6Y_Q$, and the unordered pair
$\{Y_u,Y_d\}=\{4Y_Q,-2Y_Q\}$; the requirement that the charged lepton Yukawa term
$-\bar L_L\,Y_e\,H\,e_R$ in \eqref{eq:yukawas} is gauge invariant fixes $Y_H=Y_L-Y_e=3Y_Q$, and then the quark Yukawa
relations \eqref{eq:yuk-rel} select the ordered assignment in \eqref{eq:hyper-solution} with $u_R$ the singlet that
couples via $\tilde H$.
\end{theorem}

\begin{proof}
Assume the anomaly conditions \eqref{eq:A3}--\eqref{eq:A1} with $Y_Q\neq 0$.

From \eqref{eq:A2} we obtain $Y_L=-3Y_Q$.
From \eqref{eq:A3} we obtain $Y_u+Y_d=2Y_Q$.
From \eqref{eq:Agrav}, using $Y_u+Y_d=2Y_Q$, we obtain
\[
0=6Y_Q+2Y_L-3(Y_u+Y_d)-Y_e
=6Y_Q+2(-3Y_Q)-3(2Y_Q)-Y_e
=-6Y_Q-Y_e,
\]
so $Y_e=-6Y_Q$.

Now write $Y_d=2Y_Q-Y_u$ and substitute $Y_L=-3Y_Q$, $Y_e=-6Y_Q$ into the cubic constraint \eqref{eq:A1}.
Expanding (Appendix~\ref{app:cubic-check}) yields the quadratic equation
\begin{equation}\label{eq:Yu-quad}
Y_u^2-2Y_QY_u-8Y_Q^2=0,
\end{equation}
and hence $Y_u\in\{4Y_Q,-2Y_Q\}$ with $Y_d=2Y_Q-Y_u\in\{-2Y_Q,4Y_Q\}$.
Thus the anomaly equations fix the unordered pair $\{Y_u,Y_d\}=\{4Y_Q,-2Y_Q\}$.

Now impose existence of the charged lepton Yukawa operator $-\bar L_L\,Y_e\,H\,e_R$ in \eqref{eq:yukawas}.
Gauge invariance under $\U(1)_Y$ requires the third relation in \eqref{eq:yuk-rel}, $Y_e=Y_L-Y_H$, hence
\begin{equation}\label{eq:YH-from-anom}
Y_H = Y_L-Y_e = (-3Y_Q)-(-6Y_Q)=3Y_Q.
\end{equation}
With this $Y_H$, the quark Yukawa relations \eqref{eq:yuk-rel} force
$Y_u=Y_Q+Y_H=4Y_Q$ and $Y_d=Y_Q-Y_H=-2Y_Q$, yielding \eqref{eq:hyper-solution}.

The overall normalization $Y_Q$ is redundant (absorbed into $g_1$) and the overall sign corresponds to charge
conjugation (Theorem~\ref{thm:hypercharge}).
\end{proof}

\begin{corollary}[Charge lattice rigidity]\label{cor:charge-lattice}
Within the declared load, all anomaly-free hypercharge assignments are proportional to the standard
solution; the only freedom is an overall normalization (absorbed into $g_1$) and a global sign.
Consequently, once the redundancy of hypercharge normalization is fixed (Theorem~\ref{thm:hypercharge}),
electric charge $Q=T^3+Y$ is quantized on a discrete lattice.
\end{corollary}

\begin{remark}[Neutrino mass does not affect the forcing of hypercharge]\label{rem:nu-hyper}
Neutrino oscillation bookkeeping (Section~3.5) can be accommodated either by adding gauge-singlet
right-handed neutrinos (which are anomaly-inert if $Y_{\nu_R}=0$) or by an effective Weinberg operator.
Neither option alters the hypercharge forcing in Theorem~\ref{thm:anomaly-yukawa} because the anomaly
constraints and the charged Yukawa compatibility are already fixed by the charged chiral skeleton.
\end{remark}
\subsection{Hypercharge normalization and charge quantization}

\begin{theorem}[Hypercharge normalization collapse (redundancy elimination)]\label{thm:hypercharge}
Fix the declared load and the anomaly/Yukawa skeleton.
The apparent continuous freedom to rescale hypercharge assignments is redundancy: it is removable by
admissible redescription (a simultaneous rescaling of $Y$ and $g_1$) and collapses to a canonical charge
lattice choice, unique up to overall sign.
Consequently, electric charge $Q=T^3+Y$ is quantized on a discrete lattice fixed by the skeleton.
\end{theorem}

\begin{proof}
Start in the unbroken gauge skin where hypercharge enters only through the product $g_1 Y$ in the
covariant derivative. For any nonzero real $\lambda$, the simultaneous rescaling
\begin{equation}\label{eq:rescale}
Y\mapsto \lambda Y,\qquad g_1\mapsto g_1/\lambda
\end{equation}
leaves $g_1 Y$ invariant and therefore leaves all gauge-phase interactions unchanged. Thus $\lambda$ is
\emph{not} a standing-relevant degree; it is a redundancy in the gauge-phase parameterization.

In the broken-phase skin, the unbroken electromagnetic generator is
\begin{equation}\label{eq:Q}
Q=T^3+Y,
\end{equation}
with $T^3$ fixed by the non-abelian representation content. Given Corollary~\ref{cor:charge-lattice},
all hypercharges are proportional to a single normalization choice; choosing the representative in which
the charge lattice is minimal (no common integer factor) fixes the redundancy \eqref{eq:rescale}.
Only the global sign remains (charge conjugation).

Therefore hypercharge normalization is fixed up to sign as a redundancy elimination, and the resulting
electric charges lie on a discrete lattice determined by the skeleton (hence are quantized).
\end{proof}

\subsection{Absence of tree-level FCNC (fully worked)}

\begin{theorem}[No tree-level flavor-changing neutral currents in the minimal load]\label{thm:fcnc}
Within the declared load (single Higgs doublet; generation-universal gauge charges), neutral currents
remain flavor-diagonal at tree level in the mass basis. Tree-level flavor-changing neutral currents
(FCNC) require forbidden additional structure (extra Higgs multiplets, non-universal gauge generators,
or compensating sectors).
\end{theorem}

\begin{proof}
We show the diagonalization explicitly.

\paragraph{(1) Yukawa matrices and mass diagonalization.}
For $n$ generations, the quark Yukawa terms are
\begin{equation}
\mathcal{L}_Y\supset -\bar Q_L Y_d H d_R - \bar Q_L Y_u \tilde H u_R + \mathrm{h.c.},
\end{equation}
with complex $n\times n$ matrices $Y_u,Y_d$. After EWSB, $M_{u,d}=Y_{u,d}v/\sqrt2$.
There exist unitary matrices $U_{uL},U_{uR},U_{dL},U_{dR}$ such that
\begin{equation}
U_{uL}^\dagger M_u U_{uR}=\diag(m_u,m_c,\dots),\qquad
U_{dL}^\dagger M_d U_{dR}=\diag(m_d,m_s,\dots).
\end{equation}
Define mass eigenstate fields
\begin{equation}
u_L'=U_{uL}^\dagger u_L,\quad u_R'=U_{uR}^\dagger u_R,\quad
d_L'=U_{dL}^\dagger d_L,\quad d_R'=U_{dR}^\dagger d_R.
\end{equation}

\paragraph{(2) Charged currents produce CKM.}
The left-handed charged current in the gauge basis is
\begin{equation}
\mathcal{L}_{\CC}\supset -\frac{g_2}{\sqrt2}\,W^+_\mu\,\bar u_L\gamma^\mu d_L + \mathrm{h.c.}
\end{equation}
In the mass basis,
\begin{equation}
\bar u_L\gamma^\mu d_L
=
\bar u_L' U_{uL}\gamma^\mu U_{dL}^\dagger d_L'
=
\bar u_L'\gamma^\mu (U_{uL}^\dagger U_{dL})\,d_L'.
\end{equation}
Thus the CKM matrix is $V\equiv U_{uL}^\dagger U_{dL}$ and charged currents are flavor-mixing.

\paragraph{(3) Neutral currents remain diagonal.}
The $Z$ boson couples to neutral currents of the form
\begin{equation}\label{eq:neutral-current}
\mathcal{L}_{\NC}\supset Z_\mu\,\bar f\,\gamma^\mu\left(g_L^f P_L+g_R^f P_R\right)f,
\end{equation}
where $g_{L,R}^f$ depend only on the gauge charges of $f$ (and are generation-universal in the declared load).
For example, for left-handed up-type quarks $f=u$ one has
\begin{equation}
\bar u_L\gamma^\mu u_L = \bar u_L' U_{uL}\gamma^\mu U_{uL}^\dagger u_L' = \bar u_L'\gamma^\mu u_L',
\end{equation}
since $U_{uL}$ is unitary. The same holds for $d_L$, $u_R$, $d_R$, and all leptons: a flavor-space unitary
rotation cancels in $\bar f \gamma^\mu f$. Therefore all neutral currents in \eqref{eq:neutral-current} are
flavor-diagonal at tree level in the mass basis.

\paragraph{(4) Why FCNC require extra structure.}
Tree-level FCNC arise if (and essentially only if) the neutral current couplings are not proportional to
the identity in flavor space \emph{in the same basis that diagonalizes the mass matrices}. In the declared
load, gauge charges are generation-universal and the neutral current generators are fixed, so the couplings
are proportional to identity and remain diagonal after unitary rotations.
To obtain non-universal neutral couplings one must add new Higgs multiplets with misaligned Yukawas,
new gauge generators with non-universal charges, or compensating sectors---all forbidden by the declared load.
\end{proof}

\section{Mixing and rephasing: saturation and redundancy (fully shown)}

\subsection{Parameter counting and invariants}
We state and prove the standard parameter-counting result, because it is the audit point for
``mixing saturation'' and ``phase exhaustion''.

\begin{theorem}[Physical mixing parameters as quotient invariants]\label{thm:mixing-count}
For $n$ generations of quarks in the declared load, the Yukawa sector contains, after quotienting by
admissible field redefinitions, the following physical invariants:
\begin{itemize}[leftmargin=2.2em]
\item $2n$ quark masses;
\item $\frac{n(n-1)}{2}$ mixing angles;
\item $\frac{(n-1)(n-2)}{2}$ CP-odd phases (Dirac-type).
\end{itemize}
For $n=3$, this is $6$ masses, $3$ angles, and $1$ phase.
\end{theorem}

\begin{proof}
The two Yukawa matrices $Y_u,Y_d$ are complex $n\times n$ matrices, hence contain $2n^2$ real parameters each,
for a total of $4n^2$.

Admissible field redefinitions in the quark kinetic terms permit independent unitary rotations
\begin{equation}
Q_L\mapsto U_Q Q_L,\qquad u_R\mapsto U_u u_R,\qquad d_R\mapsto U_d d_R,
\end{equation}
with $U_Q,U_u,U_d\in U(n)$, without changing gauge interactions (they are flavor-universal).
These transformations act on Yukawas by
\begin{equation}
Y_u\mapsto U_Q Y_u U_u^\dagger,\qquad
Y_d\mapsto U_Q Y_d U_d^\dagger.
\end{equation}
The group $U(n)^3$ has real dimension $3n^2$. However, one overall $U(1)$ corresponds to baryon number and
does not change $Y_u,Y_d$ (it acts trivially on both); therefore only $3n^2-1$ parameters can be removed.
Hence the physical parameter count is
\begin{equation}
4n^2-(3n^2-1)=n^2+1.
\end{equation}
On the other hand, $n^2+1$ decomposes as
\begin{equation}
n^2+1 = 2n + \frac{n(n-1)}{2} + \frac{(n-1)(n-2)}{2},
\end{equation}
corresponding to $2n$ masses, $\frac{n(n-1)}{2}$ angles, and $\frac{(n-1)(n-2)}{2}$ phases. This matches
the standard CKM parameterization count. For $n=3$, one obtains $10$ invariants: $6$ masses, $3$ angles, $1$
phase.
\end{proof}

\subsection{Saturation and redundancy statements}

\begin{corollary}[Unitary mixing saturation]\label{cor:mixing-saturation}
In the declared load, all unitary mixing freedom is exhausted by admissible basis choice and Yukawa
diagonalization. The CKM (and similarly PMNS, in the appropriate neutrino mass assumptions) is a
\emph{quotient invariant} of Yukawa operator content, not an additional independent tensor load.
\end{corollary}

\begin{remark}[Not ``unphysical mixing'']\label{rem:not-unphysical-mixing}
Corollary~\ref{cor:mixing-saturation} is often misread as a denial of physical mixing.
It is not. It states that, within the declared load, mixing data are \emph{quotient invariants} of the
Yukawa operator content (misalignment of $Y_u$ and $Y_d$) rather than an additional independent tensor
load that can be specified on top of the Yukawas without changing witness content.
\end{remark}


\begin{corollary}[Rephasing redundancy exhaustion]\label{cor:rephasing}
All rephasing freedom beyond the physical CP-odd quotient phases is redundancy. Any claim of additional
independent phases requires representational privilege or hidden structure and is therefore a standing
failure (Definition~\ref{def:failure}).
\end{corollary}

\section{Exhaustion: no further admissible equivalence completion remains}

\subsection{What ``equivalence completion'' means here}
An \emph{equivalence completion} is any forced identification of an apparent parameter freedom as redundancy
under admissible redescription and admissible transport, resulting in a smaller witness-determined tensor
load for a fixed scope.

\subsection{Admissibility envelope and a normal-form map}\label{sec:envelope-nf}
The word ``exhaustion'' is strong, so we explicitly delimit what is being exhausted.

\begin{definition}[Admissibility envelope used here]\label{def:envelope}
The \emph{admissibility envelope} $\Ad_{\Load}$ used for the exhaustion claim is the closure of the
renormalizable gauge-phase skin under finite compositions of the following \emph{explicitly licensed}
moves inside the declared load:
\begin{enumerate}[label=(\roman*),leftmargin=2.2em]
\item admissible redescriptions (Definition~\ref{def:redescription}): gauge-fixing changes, field
redefinitions, flavor-basis rotations, and rephasings;
\item broken$\leftrightarrow$unbroken transport by explicit Higgs expansion and mass diagonalization;
\item regime transport by explicit matching/integration-out inside $\Load$ (e.g.\ $W$-exchange $\to$
four-fermion, $A/Z$ rotation $\to$ QED), with an explicit scope and error order when approximations are used
(Definition~\ref{def:scope} and Remark~\ref{rem:eft});
\item on-shell reduction and high-energy limits with controlled corrections (e.g.\ the longitudinal--Goldstone
transport).
\end{enumerate}
We do \emph{not} claim to classify arbitrary skins outside $\Ad_{\Load}$; leaving the envelope is a scope
change by Definition~\ref{def:scope-preserving}.
\end{definition}

\begin{remark}[Renormalization scheme and running are scope data]\label{rem:rg-scope}
If a scope $S$ includes variation in renormalization scale or scheme, then those choices are part of the scope
specification (Definition~\ref{def:scope}) and any change of scale/scheme must be exhibited as an explicit transport
(map of couplings and counterterms) inside the declared load.
In this paper, couplings are treated as witnesses at a chosen reference scale; renormalization-group running to
another scale is a controlled transport, not an additional standing degree of freedom.
Claims that treat scheme dependence as an extra ``physical parameter'' without an explicit transport certificate
are standing failures of type (b) or (c) (Definition~\ref{def:failure}).
\end{remark}


\begin{definition}[Witness normal form for a scope]\label{def:witness-normal-form}
Fix a scope $S$.
A skin $\Skin$ is in \emph{witness normal form for $S$} if its explicit parameterization consists only of
(i) discrete-skeleton data required by $S$ (e.g.\ hypercharge class, representation labels), and
(ii) a minimal invariant witness set $\Wit_S(\Skin)$ whose values determine all $S$-predictions.
\end{definition}

\begin{definition}[Canonical standing normal form skin]\label{def:NF-skin}
Fix a scope $S$ and a skin $\Skin$ with parameter space $P(\Skin)$.
Define the \emph{canonical standing normal form} $\mathrm{NF}_S(\Skin)$ to be the bookkeeping skin with
parameter space
\[
P(\mathrm{NF}_S(\Skin))\;:=\;P(\Skin)/\!\sim_S
\]
(Definition~\ref{def:standing-quotient}), and prediction rule inherited from $\Skin$ by
\begin{equation}\label{eq:NF-pred}
\mathrm{Pred}_{\mathrm{NF}_S(\Skin)}(\mathcal{O};[\mathbf{p}])\;:=\;\mathrm{Pred}_{\Skin}(\mathcal{O};\mathbf{p})
\qquad (\mathcal{O}\in S).
\end{equation}
This is well-defined because $\mathbf{p}\sim_S\mathbf{p}'$ implies equality of all $S$-predictions.
No canonical representative is assumed: only standing classes $[\mathbf{p}]$ appear.
\end{definition}

\begin{lemma}[Compilation lemma: quotient transport and universal property]\label{lem:compile}
Let $S$ be a fixed scope and $\Skin$ any skin.
Let $q:P(\Skin)\to P(\Skin)/\!\sim_S$ denote the quotient map.
Then:
\begin{enumerate}[label=(\roman*),leftmargin=2.2em]
\item $q$ defines an admissible transport $q:\Skin\to \mathrm{NF}_S(\Skin)$ that is transport-complete for $S$.
\item (Universal property.) For any set $X$ and any map $F:P(\Skin)\to X$ that is constant on $\sim_S$-classes
(i.e.\ $\mathbf{p}\sim_S\mathbf{p}'\Rightarrow F(\mathbf{p})=F(\mathbf{p}')$), there exists a unique map
$\widetilde F:P(\Skin)/\!\sim_S\to X$ such that $F=\widetilde F\circ q$.
\end{enumerate}
\end{lemma}

\begin{proof}
For (i), admissibility is immediate: $q$ introduces no new load and preserves standing by construction.
Transport completeness holds because if $\mathbf{p},\mathbf{p}'$ yield different $S$-predictions, then
$\mathbf{p}\not\sim_S \mathbf{p}'$, hence $q(\mathbf{p})\neq q(\mathbf{p}')$; the normal form retains exactly the
standing distinctions.

For (ii), define $\widetilde F([\mathbf{p}]):=F(\mathbf{p})$. This is well-defined precisely because $F$ is constant
on $\sim_S$-classes. Uniqueness follows because $q$ is surjective.
\end{proof}

\begin{theorem}[Skin normal form completeness for a scope]\label{thm:skin-nf}
Fix a scope $S$.
\begin{enumerate}[label=(\roman*),leftmargin=2.2em]
\item (Terminal compression.) For any admissible transport $T:\Skin\to\Skin'$ that is scope-preserving for $S$,
the induced map on parameter spaces factors uniquely through the quotient transport $q:\Skin\to\mathrm{NF}_S(\Skin)$.
Equivalently: no admissible same-scope map can depend on more than the standing class $[\mathbf{p}]$.
\item (Completeness.) Within the admissibility groupoid $\Ad_S$ (Theorem~\ref{thm:coherence-groupoid}), the normal
form is a complete invariant: if $\Skin_A\simeq \Skin_B$ in $\Ad_S$, then
$\mathrm{NF}_S(\Skin_A)\cong \mathrm{NF}_S(\Skin_B)$ as standing state spaces with prediction rule.
\end{enumerate}
\end{theorem}

\begin{proof}
(i) Scope-preserving admissible transport means: if $\mathbf{p}\sim_S\mathbf{p}'$ then the transported states are
$S$-indistinguishable in $\Skin'$, hence $T(\mathbf{p})$ depends only on $[\mathbf{p}]$.
Thus the underlying map on parameters is constant on $\sim_S$-classes and factors uniquely through the quotient by
Lemma~\ref{lem:compile}(ii).

(ii) If $\Skin_A\simeq\Skin_B$ in $\Ad_S$, then by definition there exist transport-complete admissible transports
in both directions. Lemma~\ref{lem:quotient-iso} then gives a bijection of standing quotients preserving
$S$-predictions, which is exactly an isomorphism
$\mathrm{NF}_S(\Skin_A)\cong\mathrm{NF}_S(\Skin_B)$.
\end{proof}

\begin{lemma}[Existence/uniqueness of witness normal form inside $\Ad_{\Load}$]\label{lem:nf}
For each scope treated in Sections~3--5, every skin in the envelope $\Ad_{\Load}$ can be brought, by a finite
sequence of admissible redescriptions and admissible transports, to witness normal form for that scope.
Any two witness normal forms for the same scope differ only by admissible redescription.
\end{lemma}

\begin{proof}
Existence is constructive: Sections~3--5 explicitly exhibit the required transports, redundancy quotients,
and matching maps for the relevant scope classes (contact CC, pure QED, CP--odd QCD, oscillations,
high-energy longitudinal limits, and flavor-quotient mixing), thereby displaying explicit witness-coordinate
presentations inside the envelope.

Uniqueness is a quotient statement: any two witness normal forms for the same $S$ present the same standing
state space and hence differ only by admissible redescription.
Formally, the quotient $P(\Skin)/\!\sim_S$ is canonical and any standing-preserving map factors through it
uniquely (Lemma~\ref{lem:compile}); there is no additional standing left to encode.
\end{proof}

\begin{theorem}[Equivalence exhaustion (relative to $\Ad_{\Load}$)]\label{thm:exhaustion}
Within the declared load, and relative to the admissibility envelope $\Ad_{\Load}$ (Definition~\ref{def:envelope}),
there is no further admissible equivalence completion for the scope family treated in Sections~3--5 beyond:
\begin{itemize}[leftmargin=2.2em]
\item the parameter normal-form collapses (Section~3),
\item the forced discrete skeleton (Section~4),
\item the mixing/rephasing quotient classification (Section~5).
\end{itemize}
Any purported additional completion must either (i) leave $\Ad_{\Load}$ (a scope change),
(ii) introduce new load, or (iii) invoke a standing-failure mode.
\end{theorem}

\begin{proof}
Fix one of the scopes treated in Sections~3--5 and let $\Skin$ be any skin in $\Ad_{\Load}$.
By Lemma~\ref{lem:nf}, $\Skin$ reduces (by explicit licensed moves) to a witness normal form presentation for $S$
inside the envelope.
Independently of which presentation is used, the standing content for $S$ is the canonical quotient
$P(\Skin)/\!\sim_S$ (Definition~\ref{def:standing-quotient}), and $\mathrm{NF}_S(\Skin)$ is the canonical
compression (Definition~\ref{def:NF-skin}).

A further ``equivalence completion'' would have to identify (or remove) additional degrees of freedom while
preserving all $S$-predictions.
But any admissible scope-preserving transport out of $\Skin$ factors uniquely through the quotient transport
to $\mathrm{NF}_S(\Skin)$ (Theorem~\ref{thm:skin-nf}(i)).
Thus the quotient already performs the \emph{maximal} standing-preserving identification: any stricter
identification corresponds to a relation coarser than $\sim_S$, and therefore would identify two points that
are $S$-distinguishable by definition of $\sim_S$.
That changes an $S$-prediction and is therefore not a same-scope admissible completion.

Equivalently, in witness normal form every explicit parameter is either discrete skeleton data required by $S$ or
part of a \emph{minimal} invariant witness set determining $S$-predictions.
Removing a witness changes some $S$-prediction; keeping a parameter while denying its witness status contradicts
Theorem~\ref{thm:witness-counting} (or requires denying transport completeness), which is a standing failure.

Therefore no additional admissible completion exists \emph{inside} $\Ad_{\Load}$. Any escape must import new
transport moves or new structure (scope change/new load) or invoke standing-failure modes.
\end{proof}


\section{Capstone: the Standard-Model bookkeeping interior is closed (standalone statement)}

\begin{definition}[SM bookkeeping interior (in this paper)]\label{def:interior}
The \emph{\SM bookkeeping interior} is the equivalence class (under admissible redescription) of skins
within the declared load whose discrete skeleton satisfies:
\begin{enumerate}[label=(\roman*),leftmargin=2.2em]
\item anomaly/Yukawa compatibility (Theorem~\ref{thm:anomaly-yukawa});
\item hypercharge/charge-lattice rigidity (Corollary~\ref{cor:charge-lattice} and Theorem~\ref{thm:hypercharge});
\item transport-closed parameter normal form (Proposition~\ref{prop:normal-form});
\item flavor/phase quotient classification (Theorem~\ref{thm:mixing-count}).
\end{enumerate}
\end{definition}

\begin{theorem}[Interior closure and uniqueness (relative to the declared load)]\label{thm:closure}
Within the declared load, the \SM bookkeeping interior is closed and unique up to admissible redescription:
any purported same-scope ``extension'' either (a) reduces to the same interior by admissible redescription,
or (b) violates standing preservation by importing forbidden load, representational privilege, or illicit
scope transport.
\end{theorem}

\begin{proof}
The discrete skeleton is forced up to overall normalization/sign by Theorem~\ref{thm:anomaly-yukawa} and
Corollary~\ref{cor:charge-lattice}, and normalization freedom is redundancy by Theorem~\ref{thm:hypercharge}.
Within that skeleton, all apparent continuous freedoms not tied to invariant witnesses collapse under the
transport certificates of Section~3, yielding the parameter normal form (Proposition~\ref{prop:normal-form}).
Mixing and phase freedoms are fully classified as quotient invariants (Theorem~\ref{thm:mixing-count}),
and tree-level FCNC are excluded in the declared load (Theorem~\ref{thm:fcnc}).
Finally, Theorem~\ref{thm:exhaustion} closes the equivalence space: no further admissible equivalence
completion exists.

Therefore any proposed modification that stays inside the declared load is either an admissible redescription
of existing structure (hence reduces), or it introduces new operators/fields/sectors or privileges a
representation to maintain standing---which is inadmissible by Definition~\ref{def:failure} and
Remark~\ref{rem:fail-closed}. Uniqueness is thus relative to the declared load and admissibility envelope.
\end{proof}

\section{Certification: irreducible inputs and structural non-derivability}\label{sec:certification}

\subsection{The full renormalizable scope and the witness ledger}

The collapses in Sections~3--5 are \emph{scope-indexed}: in restricted regimes (contact CC, pure QED, etc.)
some couplings collapse to fewer witnesses. A hostile reader may then ask which data remain standing-relevant
once the scope is enlarged back to the full renormalizable bookkeeping interior.

\begin{definition}[Full renormalizable \SM scope]\label{def:scope-SM-ren}
Let $S_{\SM}^{\mathrm{ren}}$ be the scope consisting of all predictions of the most general
\emph{renormalizable} Lagrangian built from the declared load (Definition~\ref{def:declared-load}), including
processes sensitive to $W$, $Z$, Higgs exchange, QCD interactions, and flavor-changing charged currents.
No low-energy truncation is assumed; basis/gauge choices are treated as admissible redescriptions.
\end{definition}

\begin{proposition}[Witness normal form for $S_{\SM}^{\mathrm{ren}}$]\label{prop:SM-ren-NF}
Fix the declared load and the full renormalizable scope $S_{\SM}^{\mathrm{ren}}$.
Up to admissible redescription (gauge fixing, field redefinitions, flavor rotations, and rephasings),
the standing-relevant data reduce to the following witness ledger:
\begin{enumerate}[label=(\roman*),leftmargin=2.2em]
\item \textbf{Discrete skeleton:} gauge group and representation content (Definition~\ref{def:declared-load}),
generation count $n_g$, and the hypercharge class fixed up to overall sign (Theorem~\ref{thm:anomaly-yukawa}
and Corollary~\ref{cor:charge-lattice}).
\item \textbf{Gauge-sector witnesses:} three gauge couplings $(g_s,g_2,g_1)$ at a chosen renormalization
scheme/scale (or any equivalent invariant reparameterization, e.g.\ $(\alpha_s, e,\sin^2\theta_W)$ once the
scope includes $Z$ exchange).
\item \textbf{Higgs-sector witnesses:} two independent invariants of the renormalizable Higgs potential, e.g.\
$(v,m_h)$ (equivalently $(\mu^2,\lambda)$ in $V(H)=-\mu^2 H^\dagger H+\lambda(H^\dagger H)^2$).
\item \textbf{Yukawa/flavor witnesses:} fermion mass eigenvalues (equivalently Yukawa singular values) together
with the quotient mixing invariants (angles and CP phases, or equivalently Jarlskog-type invariants) as
classified in Section~5.
\item \textbf{CP--odd strong-sector witness:} the unique invariant $\bar\theta$ (Lemma~\ref{lem:theta}).
\item \textbf{Neutrino-sector witnesses (if neutrino masses are included in scope):} mass-squared differences and
rephasing-invariant mixing products as in Lemma~\ref{lem:osc}, with the precise ledger depending on the
Dirac/Majorana load choice (a scope/load choice, not fixed here).
\end{enumerate}
No other continuous parameter in the renormalizable Lagrangian has independent standing for
$S_{\SM}^{\mathrm{ren}}$.
\end{proposition}

\begin{proof}
Within the declared load, the most general renormalizable Lagrangian consists of:
\begin{enumerate}[label=(\alph*),leftmargin=2.2em]
\item gauge kinetic terms for $\SU(3)_c$, $\SU(2)_L$, and $\U(1)_Y$ with couplings $(g_s,g_2,g_1)$;
\item fermion kinetic terms with covariant derivatives fixed by the representation skeleton (no new parameters);
\item the Higgs kinetic term and a renormalizable potential $V(H)=-\mu^2 H^\dagger H+\lambda(H^\dagger H)^2$;
\item the Yukawa operator family \eqref{eq:yukawas} with complex matrices $(Y_u,Y_d,Y_e)$;
\item the CP--odd topological density term with coefficient $\theta$, whose standing-relevant invariant is
$\bar\theta$ by Lemma~\ref{lem:theta}.
\end{enumerate}
There are no additional renormalizable gauge-invariant operators built from the declared fields.

Admissible redescriptions remove gauge-fixing parameters, field-redefinition artifacts, and the flavor-basis
redundancies of Section~5. The hypercharge class is forced up to normalization/sign by
Theorem~\ref{thm:anomaly-yukawa} and Corollary~\ref{cor:charge-lattice}, and the normalization freedom is
redundancy by Theorem~\ref{thm:hypercharge}. The remaining independent invariant witnesses are exactly the items
listed in (i)--(vi).
\end{proof}

\begin{remark}[Why this does not contradict the low-energy collapses]\label{rem:global-vs-restricted}
Proposition~\ref{prop:SM-ren-NF} concerns the full renormalizable scope $S_{\SM}^{\mathrm{ren}}$.
In restricted scopes (e.g.\ $S_{\CC}^{(6)}$ or $S_{\EM}^{(\gamma)}$), transport collapses reduce witness rank,
and fewer combinations (e.g.\ $G_F$ or $e$) suffice. Enlarging the scope reintroduces additional witnesses
(e.g.\ $m_W,m_Z,\sin^2\theta_W$) exactly as in Remarks~\ref{rem:cc-scope} and \ref{rem:em-scope}.
\end{remark}

\subsection{Relational identities vs.\ ``ontic'' value derivations}

\begin{definition}[Relational identity vs.\ internal value derivation]\label{def:ontic}
A \emph{relational identity} is a standing-preserving relation among invariant witnesses that holds across all
admissible skins in a scope (e.g.\ $m_W=g_2 v/2$, $G_F=1/(\sqrt2\,v^2)$, or $\bar\theta=\theta+\arg\det M$).
An \emph{internal value derivation} claim is the stronger assertion that a specific numerical value of a witness
parameter (e.g.\ $\alpha_s(m_Z)$, a Yukawa eigenvalue, or a mixing angle) is forced \emph{purely} by interior
constraints without importing any additional standing.
\end{definition}

\begin{theorem}[Structural non-derivability of numerical witness values]\label{thm:non-derivability}
Within the declared load and the admissibility discipline of this paper, no internal value derivation of an
irreducible witness parameter in Proposition~\ref{prop:SM-ren-NF} is admissible.
Any such claim necessarily invokes at least one standing-failure mode
(Definition~\ref{def:failure})---typically representational privilege, illicit scope transport (e.g.\ measure or
typicality postulates), or compensating load.
\end{theorem}

\begin{proof}
Let $w$ be an irreducible invariant witness for $S_{\SM}^{\mathrm{ren}}$.
By definition, $w$ is invariant under admissible redescriptions, but its \emph{value} is not fixed by the
relations among witnesses exhibited inside the declared load: otherwise it would not be irreducible.

An internal value derivation would have to supply additional structure that selects $w=w_0$ without reference to
external boundary data. Any such selection must come from somewhere:
\begin{itemize}[leftmargin=2.2em]
\item If it is asserted as necessary by privileging a particular representation, parameterization, or basis as
identity-grounding, then the claim violates redescription stability and is a representational-privilege failure.
\item If it is asserted as ``generic''/``typical''/``measure-selected'', then the required measure/typicality
operator is not generated by standing-preserving transport inside the declared load; importing it is illicit
scope transport.
\item If it is produced by adding new fields/operators/sectors or an external ``selector'' principle, that is
compensating load forbidden by Definition~\ref{def:declared-load}.
\end{itemize}
These exhaust the ways to force a numerical value without empirical boundary input. Therefore any purported
internal value derivation violates admissibility by invoking a standing-failure mode.
\end{proof}

\begin{remark}[Classification, not a ban]\label{rem:nonfund-class}
Theorem~\ref{thm:non-derivability} is a classification statement: it does not forbid scope changes (e.g.\ UV
completions) that introduce additional standing-relevant witnesses. It states only that the numerical values
of the irreducible witnesses of the \SM interior are not derivable \emph{inside} that interior without breaking
the admissibility discipline.
\end{remark}

\section{Hostile-referee attack map (standalone)}

The table below lists common attack moves and the exact result that must be denied to reopen the space.
Any objection not listed must identify which hypothesis of the cited result is being rejected.

\begin{center}
\small
\begin{tabular}{@{}p{0.56\textwidth}p{0.36\textwidth}@{}}
\toprule
\textbf{Objection / escape attempt} & \textbf{Must deny or else blocked by} \\
\midrule
``You are conflating redundancy with regime suppression / integrating-out. The witness theorem is not justified.'' &
Definitions~\ref{def:scope}--\ref{def:standing-quotient}, Lemma~\ref{lem:quotient-iso},
Remark~\ref{rem:eft}, and Theorem~\ref{thm:witness-counting}.\\
``Declared load is a hidden anchor: you assumed the SM field content.'' &
Remark~\ref{rem:not-derive-gauge}, Section~\ref{sec:load-envelope}, and Proposition~\ref{prop:added-load}.\\
``You did not derive the gauge group / representation skeleton, so this is not a derivation of the SM.'' &
Remark~\ref{rem:not-derive-gauge} (typed scope: closure relative to a fixed load).\\
``There may exist other admissible skins or further equivalences you did not consider.'' &
Remark~\ref{rem:no-existence-by-assertion} and Definition~\ref{def:envelope} (construction requirement).\\
``$G_F$ matching is just a convention; $g_2$ and $v$ are independent in CC physics.'' &
Lemma~\ref{lem:cc} (explicit matching \eqref{eq:GFmatch}) \emph{for scope $S_{\CC}^{(6)}$} and
Remark~\ref{rem:cc-scope} (scope warning).\\
``$e$ is not the only electromagnetic witness; $g_1$ and $g_2$ are separately physical in EM.'' &
Lemma~\ref{lem:em} (rotation \eqref{eq:rotAZ} and identification \eqref{eq:e-def}) \emph{for scope
$S_{\EM}^{(\gamma)}$} and Remark~\ref{rem:em-scope}.\\
``$\theta$ and $\arg\det M$ are independent physical parameters.'' &
Lemma~\ref{lem:theta} (explicit shifts \eqref{eq:argdet-shift}--\eqref{eq:theta-shift}).\\
``You overstate anomaly$\to$Yukawa forcing / the algebra is inconsistent.'' &
Section~\ref{sec:anom-yuk}, Proposition~\ref{prop:constraint-coupling}, and
Theorem~\ref{thm:anomaly-yukawa} (fully worked, convention-consistent forcing).\\
``Hypercharge normalization is a genuine extra continuous freedom.'' &
Theorem~\ref{thm:hypercharge} (rescaling redundancy \eqref{eq:rescale}).\\
``Mixing saturation denies physical mixing.'' &
Theorem~\ref{thm:mixing-count} and Remark~\ref{rem:not-unphysical-mixing}.\\
``Tree-level FCNC can appear without extra structure; your theorem is not new.'' &
Theorem~\ref{thm:fcnc} (explicit unitary cancellation; novelty is the scope/standing interpretation).\\
``You claim to determine the numerical values of the remaining \SM parameters.'' &
Definition~\ref{def:ontic}, Theorem~\ref{thm:non-derivability}, and Remark~\ref{rem:nonfund-class}.\\
``There are further equivalence completions left; your exhaustion claim is hand-wavy.'' &
Remark~\ref{rem:no-existence-by-assertion}, Definition~\ref{def:envelope}, Lemma~\ref{lem:nf}, and
Theorem~\ref{thm:exhaustion}.\\
\bottomrule
\end{tabular}
\end{center}


\appendix

\section{Cubic anomaly check algebra}\label{app:cubic-check}
This appendix records the explicit expansion used in Theorem~\ref{thm:anomaly-yukawa} to reduce the
cubic anomaly constraint to the quadratic equation \eqref{eq:Yu-quad}.

Start from the one-generation anomaly conditions \eqref{eq:A3}--\eqref{eq:A1}.
From \eqref{eq:A2} and \eqref{eq:Agrav} we have
\[
Y_L=-3Y_Q,\qquad Y_e=-6Y_Q.
\]
From \eqref{eq:A3} we have $Y_d=2Y_Q-Y_u$.
Insert these into \eqref{eq:A1}:
\begin{align}
0
&=6Y_Q^3+2(-3Y_Q)^3-3Y_u^3-3(2Y_Q-Y_u)^3-(-6Y_Q)^3\\
&=6Y_Q^3-54Y_Q^3-3Y_u^3-3\Big(8Y_Q^3-12Y_Q^2Y_u+6Y_QY_u^2-Y_u^3\Big)+216Y_Q^3\\
&=(6-54-24+216)Y_Q^3+36Y_Q^2Y_u-18Y_QY_u^2\\
&=144Y_Q^3+36Y_Q^2Y_u-18Y_QY_u^2.
\end{align}
Assuming $Y_Q\neq 0$ and dividing by $-18Y_Q$ gives
\[
Y_u^2-2Y_QY_u-8Y_Q^2=0,
\]
which is \eqref{eq:Yu-quad}. Its solutions are $Y_u=4Y_Q$ or $Y_u=-2Y_Q$, with $Y_d=2Y_Q-Y_u$.
\section{Fujikawa derivation of the chiral Jacobian and the $\theta$ shift}\label{app:fujikawa}
The path-integral measure computation follows the standard Fujikawa method \cite{Fujikawa1979,Fujikawa1980}.

This appendix gives a line-by-line computation of the Jacobian of the fermionic path-integral measure
under a global axial rotation, and the induced shift of the $\theta$ term used in
Lemma~\ref{lem:theta}. We write the calculation in Euclidean signature (where the regulator is
well-defined) and then translate the result back to the Minkowski $\theta\,G\tilde G$ normalization.

\subsection{Setup: mode expansion and the measure}
Consider one Dirac fermion in a background $\SU(3)$ gauge field $A_\mu = A_\mu^a T^a$ with
$T^a$ Hermitian and $\Tr(T^a T^b)=\frac12\delta^{ab}$.
Let
\begin{equation}
\not{D}\equiv \gamma^\mu(\partial_\mu - i g_s A_\mu),\qquad
S_E=\int d^4x\;\bar\psi\,\not{D}\,\psi+\cdots .
\end{equation}
Choose a complete orthonormal eigenbasis of the Euclidean Dirac operator,
\begin{equation}
\not{D}\,\phi_n(x)=\lambda_n\,\phi_n(x),\qquad
\int d^4x\;\phi_n^\dagger(x)\phi_m(x)=\delta_{nm}.
\end{equation}
Expand the Grassmann fields in this basis:
\begin{equation}
\psi(x)=\sum_n a_n \phi_n(x),\qquad \bar\psi(x)=\sum_n \bar a_n \phi_n^\dagger(x),
\end{equation}
so that (up to an overall constant) the functional measure is
\begin{equation}
\mathcal{D}\psi\,\mathcal{D}\bar\psi=\prod_n da_n\,d\bar a_n.
\end{equation}

\subsection{Axial rotation and the formal Jacobian}
Perform a global axial rotation with parameter $\alpha$,
\begin{equation}\label{eq:fuj-axial}
\psi \mapsto \psi' = e^{i\alpha\gamma_5}\psi,\qquad
\bar\psi \mapsto \bar\psi'=\bar\psi\,e^{i\alpha\gamma_5},
\end{equation}
(where in Euclidean signature we take $\gamma_5^\dagger=\gamma_5$, $\gamma_5^2=\mathbf{1}$).
At the level of mode coefficients,
\begin{equation}
a_n'=\sum_m \left(\int d^4x\;\phi_n^\dagger e^{i\alpha\gamma_5}\phi_m\right)a_m
\equiv \sum_m (U_\alpha)_{nm}\,a_m,
\end{equation}
and similarly $\bar a_n'=\sum_m \bar a_m (U_\alpha)_{mn}$.
Hence
\begin{equation}
\prod_n da_n' = (\det U_\alpha)^{-1}\prod_n da_n,\qquad
\prod_n d\bar a_n' = (\det U_\alpha)^{-1}\prod_n d\bar a_n,
\end{equation}
so the full Jacobian is
\begin{equation}\label{eq:fuj-Jformal}
\mathcal{D}\psi'\,\mathcal{D}\bar\psi' = J(\alpha)\,\mathcal{D}\psi\,\mathcal{D}\bar\psi,
\qquad
J(\alpha)=(\det U_\alpha)^{-2}.
\end{equation}
For infinitesimal $\alpha$, $U_\alpha = \mathbf{1}+ i\alpha\,M + \cO(\alpha^2)$ with
\begin{equation}
M_{nm}=\int d^4x\;\phi_n^\dagger(x)\gamma_5\phi_m(x),
\end{equation}
so
\begin{equation}\label{eq:fuj-Jinf}
\ln J(\alpha)= -2\Tr\ln U_\alpha
= -2\,\Tr\!\left(i\alpha\,M\right)+\cO(\alpha^2)
= -2i\alpha \sum_n \int d^4x\;\phi_n^\dagger(x)\gamma_5\phi_n(x)+\cO(\alpha^2).
\end{equation}
The trace $\sum_n \phi_n^\dagger\gamma_5\phi_n$ is ultraviolet divergent and must be regulated.

\subsection{Gauge-covariant regulator and reduction to a local trace}
Introduce the gauge-covariant heat-kernel regulator with mass scale $M$:
\begin{equation}\label{eq:fuj-reg}
\sum_n \phi_n^\dagger(x)\gamma_5\phi_n(x)
\;\mapsto\;
\lim_{M\to\infty}\sum_n \phi_n^\dagger(x)\gamma_5 e^{-\lambda_n^2/M^2}\phi_n(x)
=
\lim_{M\to\infty}\Tr\!\left[\gamma_5 e^{-\not{D}^2/M^2}\right]_{x,x}.
\end{equation}
Define the regulated local quantity
\begin{equation}\label{eq:fuj-Ax}
\mathcal{A}(x)\equiv \lim_{M\to\infty}\Tr\!\left[\gamma_5 e^{-\not{D}^2/M^2}\right]_{x,x}.
\end{equation}
Then, to leading order in $\alpha$,
\begin{equation}\label{eq:fuj-J-A}
\ln J(\alpha) = -2i\alpha \int d^4x\;\mathcal{A}(x)
\qquad\Rightarrow\qquad
J(\alpha)=\exp\!\left(-2i\alpha\int d^4x\;\mathcal{A}(x)\right).
\end{equation}

\subsection{Compute $\not{D}^2$ explicitly}
We now reduce $\not{D}^2$ to a Laplacian plus a field-strength term. Start from
\begin{equation}
\not{D}^2=\gamma^\mu\gamma^\nu D_\mu D_\nu,
\qquad D_\mu=\partial_\mu - i g_s A_\mu.
\end{equation}
Use the decomposition
\begin{equation}
\gamma^\mu\gamma^\nu=\frac12\{\gamma^\mu,\gamma^\nu\}+\frac12[\gamma^\mu,\gamma^\nu]
= g^{\mu\nu}\mathbf{1} + \sigma^{\mu\nu},
\qquad
\sigma^{\mu\nu}\equiv \frac12[\gamma^\mu,\gamma^\nu],
\end{equation}
(where in Euclidean signature $g^{\mu\nu}=\delta^{\mu\nu}$; the calculation below is signature-agnostic
as long as the standard Clifford algebra holds).
Then
\begin{equation}
\not{D}^2 = D^\mu D_\mu + \sigma^{\mu\nu}D_\mu D_\nu.
\end{equation}
Split $D_\mu D_\nu$ into symmetric and antisymmetric parts:
\begin{equation}
D_\mu D_\nu = \frac12\{D_\mu,D_\nu\}+\frac12[D_\mu,D_\nu].
\end{equation}
Because $\sigma^{\mu\nu}$ is antisymmetric, $\sigma^{\mu\nu}\{D_\mu,D_\nu\}=0$, so
\begin{equation}
\sigma^{\mu\nu}D_\mu D_\nu = \frac12\sigma^{\mu\nu}[D_\mu,D_\nu].
\end{equation}
With our convention $D_\mu=\partial_\mu-i g_s A_\mu$, the commutator is
\begin{equation}
[D_\mu,D_\nu]= -i g_s\,G_{\mu\nu},\qquad
G_{\mu\nu} \equiv \partial_\mu A_\nu-\partial_\nu A_\mu - i g_s[A_\mu,A_\nu] = G_{\mu\nu}^a T^a.
\end{equation}
Therefore
\begin{equation}\label{eq:fuj-Dslash2}
\not{D}^2 = D^2 - \frac{i g_s}{2}\sigma^{\mu\nu}G_{\mu\nu}.
\end{equation}

\subsection{Heat-kernel expansion: isolate the $\gamma_5$ term}
Insert \eqref{eq:fuj-Dslash2} into \eqref{eq:fuj-Ax}:
\begin{equation}
\mathcal{A}(x)=\lim_{M\to\infty}\Tr\!\left[\gamma_5\exp\!\left(-\frac{1}{M^2}\left(D^2 - \frac{i g_s}{2}\sigma^{\mu\nu}G_{\mu\nu}\right)\right)\right]_{x,x}.
\end{equation}
Expand the exponential in powers of $1/M$ and keep only terms that can contribute to a $\gamma_5$
trace. Write
\begin{equation}
\exp\!\left(-\frac{1}{M^2}(D^2+X)\right)
=
e^{-D^2/M^2}\left[\mathbf{1}-\frac{X}{M^2}+\frac{1}{2}\frac{X^2}{M^4}+\cO\!\left(\frac{1}{M^6}\right)\right],
\qquad
X\equiv +\frac{i g_s}{2}\sigma^{\mu\nu}G_{\mu\nu}.
\end{equation}
Then
\begin{align}
\mathcal{A}(x)
&=\lim_{M\to\infty}\Tr\!\left[\gamma_5 e^{-D^2/M^2}\right]_{x,x} \notag\\
&\quad -\lim_{M\to\infty}\frac{1}{M^2}\Tr\!\left[\gamma_5 e^{-D^2/M^2}X\right]_{x,x} \notag\\
&\quad +\lim_{M\to\infty}\frac{1}{2M^4}\Tr\!\left[\gamma_5 e^{-D^2/M^2}X^2\right]_{x,x}
+\cdots . \notag
\end{align}
The first term vanishes because $\Tr_{\mathrm{spin}}(\gamma_5)=0$.
The second term vanishes because $\Tr_{\mathrm{spin}}(\gamma_5\sigma^{\mu\nu})=0$.
Thus the leading contribution comes from the $X^2$ term:
\begin{equation}\label{eq:fuj-leading}
\mathcal{A}(x)=\lim_{M\to\infty}\frac{1}{2M^4}\Tr\!\left[\gamma_5 e^{-D^2/M^2}\left(\frac{i g_s}{2}\sigma^{\mu\nu}G_{\mu\nu}\right)^2\right]_{x,x}.
\end{equation}

\subsection{Coincident-point kernel and the momentum integral}
To extract the leading $M\to\infty$ limit, use that the coincident-point kernel of $e^{-D^2/M^2}$
is dominated by short distances. To the order needed in \eqref{eq:fuj-leading}, one may replace
$D^2$ by the free Laplacian in the coincident kernel, because any insertion of gauge fields from $D^2$
brings additional powers of $1/M$ that vanish in the $M\to\infty$ limit.
Thus,
\begin{equation}\label{eq:fuj-kernel}
\left\langle x\middle| e^{-D^2/M^2}\middle|x\right\rangle
=
\int\frac{d^4k}{(2\pi)^4}\,e^{-k^2/M^2} + \cO\!\left(\frac{1}{M^2}\right)
=
\frac{M^4}{16\pi^2}+\cO(M^2).
\end{equation}
(We used $\int d^4k\,e^{-k^2/M^2}=\pi^2 M^4$ and divided by $(2\pi)^4$.)

Insert \eqref{eq:fuj-kernel} into \eqref{eq:fuj-leading} and pull out the slowly varying fields:
\begin{equation}\label{eq:fuj-pullout}
\mathcal{A}(x)=\lim_{M\to\infty}\frac{1}{2M^4}\left(\frac{M^4}{16\pi^2}\right)
\Tr_{\mathrm{spin,color}}\!\left[\gamma_5\left(\frac{i g_s}{2}\sigma^{\mu\nu}G_{\mu\nu}\right)^2\right] + \cO\!\left(\frac{1}{M^2}\right).
\end{equation}

\subsection{Spinor and color traces}
Compute the trace in \eqref{eq:fuj-pullout}.
First expand
\begin{equation}
\left(\frac{i g_s}{2}\sigma^{\mu\nu}G_{\mu\nu}\right)^2
=
-\frac{g_s^2}{4}\,\sigma^{\mu\nu}\sigma^{\rho\sigma}\,G_{\mu\nu}^a G_{\rho\sigma}^b\,T^a T^b.
\end{equation}
Therefore
\begin{equation}
\Tr_{\mathrm{spin,color}}\!\left[\gamma_5\left(\frac{i g_s}{2}\sigma^{\mu\nu}G_{\mu\nu}\right)^2\right]
=
-\frac{g_s^2}{4}\,\Tr_{\mathrm{spin}}\!\left(\gamma_5\sigma^{\mu\nu}\sigma^{\rho\sigma}\right)
\,\Tr_{\mathrm{color}}(T^a T^b)\,G_{\mu\nu}^a G_{\rho\sigma}^b.
\end{equation}
Using $\Tr_{\mathrm{color}}(T^a T^b)=\tfrac12\delta^{ab}$, this becomes
\begin{equation}\label{eq:fuj-tr1}
\Tr_{\mathrm{spin,color}}[\cdots]
=
-\frac{g_s^2}{8}\,\Tr_{\mathrm{spin}}\!\left(\gamma_5\sigma^{\mu\nu}\sigma^{\rho\sigma}\right)
\,G_{\mu\nu}^a G_{\rho\sigma}^a.
\end{equation}

Now compute the spinor trace. Using $\sigma^{\mu\nu}=\tfrac12[\gamma^\mu,\gamma^\nu]$ and cyclicity,
\begin{align}
\Tr_{\mathrm{spin}}\!\left(\gamma_5\sigma^{\mu\nu}\sigma^{\rho\sigma}\right)
&=\frac14 \Tr_{\mathrm{spin}}\!\left(\gamma_5[\gamma^\mu,\gamma^\nu][\gamma^\rho,\gamma^\sigma]\right)\nonumber\\
&=\frac14\Tr_{\mathrm{spin}}\!\left(\gamma_5(\gamma^\mu\gamma^\nu-\gamma^\nu\gamma^\mu)(\gamma^\rho\gamma^\sigma-\gamma^\sigma\gamma^\rho)\right)\nonumber\\
&=\frac14\Big(
\Tr(\gamma_5\gamma^\mu\gamma^\nu\gamma^\rho\gamma^\sigma)
-\Tr(\gamma_5\gamma^\mu\gamma^\nu\gamma^\sigma\gamma^\rho)\nonumber\\
&\qquad\qquad
-\Tr(\gamma_5\gamma^\nu\gamma^\mu\gamma^\rho\gamma^\sigma)
+\Tr(\gamma_5\gamma^\nu\gamma^\mu\gamma^\sigma\gamma^\rho)
\Big).
\end{align}
Use the standard identity (valid in four dimensions with $\epsilon^{0123}=+1$ in Minkowski signature,
and with corresponding Euclidean continuation)
\begin{equation}\label{eq:fuj-tr5}
\Tr(\gamma_5\gamma^\mu\gamma^\nu\gamma^\rho\gamma^\sigma)=4\,\epsilon^{\mu\nu\rho\sigma}.
\end{equation}
Because $\epsilon^{\mu\nu\sigma\rho}=-\epsilon^{\mu\nu\rho\sigma}$ and
$\epsilon^{\nu\mu\rho\sigma}=-\epsilon^{\mu\nu\rho\sigma}$, each of the four terms contributes the same
$\epsilon^{\mu\nu\rho\sigma}$. Therefore
\begin{equation}\label{eq:fuj-spintrace}
\Tr_{\mathrm{spin}}\!\left(\gamma_5\sigma^{\mu\nu}\sigma^{\rho\sigma}\right)
=
4\,\epsilon^{\mu\nu\rho\sigma}.
\end{equation}
(If one uses a different $\gamma_5$ convention, the overall sign is absorbed into the definition of
$\tilde G^{\mu\nu}$ and the sign of $\theta$; the physical invariant $\bar\theta$ in Lemma~\ref{lem:theta}
is unaffected.)

Insert \eqref{eq:fuj-spintrace} into \eqref{eq:fuj-tr1}:
\begin{equation}\label{eq:fuj-tr2}
\Tr_{\mathrm{spin,color}}[\cdots]
=
-\frac{g_s^2}{8}\,(4\epsilon^{\mu\nu\rho\sigma})\,G_{\mu\nu}^a G_{\rho\sigma}^a
=
-\frac{g_s^2}{2}\,\epsilon^{\mu\nu\rho\sigma}G_{\mu\nu}^a G_{\rho\sigma}^a.
\end{equation}
Using $\tilde G^{a\mu\nu}\equiv \tfrac12\epsilon^{\mu\nu\rho\sigma}G_{\rho\sigma}^a$, we have
$\epsilon^{\mu\nu\rho\sigma}G_{\mu\nu}^a G_{\rho\sigma}^a = 2\,G_{\mu\nu}^a\tilde G^{a\mu\nu}$.
Thus
\begin{equation}\label{eq:fuj-tr3}
\Tr_{\mathrm{spin,color}}[\cdots]
=
- g_s^2\,G_{\mu\nu}^a\tilde G^{a\mu\nu}.
\end{equation}

\subsection{Assemble the coefficient}
Combine \eqref{eq:fuj-pullout} and \eqref{eq:fuj-tr3}:
\begin{equation}
\mathcal{A}(x)
=
\frac{1}{2}\cdot\frac{1}{16\pi^2}\cdot\Big(- g_s^2\,G_{\mu\nu}^a\tilde G^{a\mu\nu}\Big)
=
-\frac{g_s^2}{32\pi^2}\,G_{\mu\nu}^a\tilde G^{a\mu\nu}.
\end{equation}
For $N_f$ flavors rotated simultaneously, $\mathcal{A}(x)$ is multiplied by $N_f$.

Insert this into \eqref{eq:fuj-J-A}:
\begin{equation}\label{eq:fuj-Jfinal}
J(\alpha)=\exp\!\left(+2i\alpha N_f\int d^4x\;\frac{g_s^2}{32\pi^2}G_{\mu\nu}^a\tilde G^{a\mu\nu}\right),
\end{equation}
where the sign follows our conventions above.

\subsection{Read off the $\theta$ shift}
The Jacobian contributes to the effective action as $\Delta S_E = -\ln J(\alpha)$, i.e.
\begin{equation}
\Delta S_E
=
-2i\alpha N_f\int d^4x\;\frac{g_s^2}{32\pi^2}G_{\mu\nu}^a\tilde G^{a\mu\nu}.
\end{equation}
Translating back to Minkowski signature and comparing with
\begin{equation}
\mathcal{L}_\theta = \frac{g_s^2}{32\pi^2}\theta\,G_{\mu\nu}^a\tilde G^{a\mu\nu},
\end{equation}
we obtain the shift
\begin{equation}
\theta \mapsto \theta + 2N_f\alpha,
\end{equation}
which is the step used in Lemma~\ref{lem:theta}.

\section{Broken-gauge Ward/Slavnov--Taylor identity and longitudinal--Goldstone equivalence}\label{app:LG-proof}
The equivalence theorem is standard; see e.g.\ \cite{Cornwall1974,Chanowitz1985}.

This appendix gives a self-contained derivation of the identity used in Lemma~\ref{lem:LG}.
We work in an $R_\xi$ gauge and use BRST invariance to derive a Slavnov--Taylor relation that links
amplitudes with an external longitudinal massive gauge boson to amplitudes with the corresponding
would-be Goldstone boson. We then take the high-energy limit $E\gg m_W$ to obtain the equivalence
of longitudinal and Goldstone amplitude classes.

\subsection{Gauge fixing with an auxiliary field}
Consider (for definiteness) the $\SU(2)$ part of the electroweak sector after spontaneous symmetry
breaking, with massive gauge bosons $W_\mu^a$ and would-be Goldstones $\pi^a$ (the argument applies
componentwise to $W^\pm$ and $Z$).
Choose the $R_\xi$ gauge-fixing function
\begin{equation}\label{eq:LG-F}
F^a \equiv \partial^\mu W_\mu^a + \xi\,m_W\,\pi^a.
\end{equation}
Introduce the Nakanishi--Lautrup auxiliary field $B^a$ and write the gauge-fixing Lagrangian as
\begin{equation}\label{eq:LG-LgfB}
\mathcal{L}_{\mathrm{gf}} = B^a F^a + \frac{\xi}{2} B^a B^a.
\end{equation}
Eliminating $B^a$ by its algebraic equation of motion gives $B^a = -\frac{1}{\xi}F^a$ and hence
\begin{equation}\label{eq:LG-Lgf}
\mathcal{L}_{\mathrm{gf}} = -\frac{1}{2\xi}(F^a)^2
= -\frac{1}{2\xi}\left(\partial\!\cdot\! W^a + \xi m_W \pi^a\right)^2,
\end{equation}
which is the standard $R_\xi$ gauge fixing used in the main text.

\subsection{Ghost term and BRST transformations}
Introduce Faddeev--Popov ghosts $c^a$ and antighosts $\bar c^a$.
The gauge-fixed action is made BRST invariant by adding the ghost term
\begin{equation}\label{eq:LG-Lgh}
\mathcal{L}_{\mathrm{gh}} = -\bar c^a\,\frac{\delta F^a}{\delta\theta^b}\,c^b,
\end{equation}
where $\theta^b$ are the infinitesimal gauge parameters and $\delta F^a/\delta\theta^b$ is computed
from the gauge transformations of $W_\mu^a$ and $\pi^a$.

BRST transformations $s$ are defined so that gauge transformations are reproduced with ghost parameters
and $s^2=0$. The only pieces we will need explicitly are the linearized transformations around the
vacuum (the nonlinear terms are fixed by the gauge algebra but do not enter the identity we derive):
\begin{align}
s\,W_\mu^a &= \partial_\mu c^a + \cdots ,\label{eq:LG-sW}\\
s\,\pi^a &= -m_W\,c^a + \cdots ,\label{eq:LG-spi}\\
s\,c^a &= -\frac{g_2}{2}\,\epsilon^{abc}c^b c^c ,\label{eq:LG-sc}\\
s\,\bar c^a &= B^a,\qquad s\,B^a=0.\label{eq:LG-sbarc}
\end{align}
A key structural fact is that the gauge-fixing plus ghost Lagrangian is BRST-exact:
\begin{equation}\label{eq:LG-exact}
\mathcal{L}_{\mathrm{gf}}+\mathcal{L}_{\mathrm{gh}} = s\!\left[\bar c^a\left(F^a + \frac{\xi}{2}B^a\right)\right].
\end{equation}
BRST invariance of the full gauge-fixed action follows immediately from $s^2=0$.

\subsection{$F^a$ is BRST-exact on shell}
Because $s\,\bar c^a = B^a$ and $B^a$ is algebraic, the $B$ equation of motion gives
\begin{equation}\label{eq:LG-BEOM}
B^a = -\frac{1}{\xi}F^a.
\end{equation}
Therefore
\begin{equation}\label{eq:LG-Fexact}
F^a = -\xi\,B^a = -\xi\,s\,\bar c^a.
\end{equation}
Thus $F^a$ is BRST-exact (after eliminating the auxiliary field). This is the algebraic origin of the
Slavnov--Taylor identity relating $\partial\!\cdot\!W^a$ and $\pi^a$.

\subsection{Vanishing of BRST-exact insertions between physical states}
Physical asymptotic states are taken to be BRST-closed (they are annihilated by the BRST charge)
and BRST-exact states are regarded as null. Consequently, matrix elements of BRST-exact operators
between physical states vanish:
\begin{equation}\label{eq:LG-BRSTnull}
\bra{\mathrm{phys}}\,s(\mathcal{O})\,\ket{\mathrm{phys}}=0.
\end{equation}
Applying \eqref{eq:LG-BRSTnull} to $\mathcal{O}=\bar c^a$ and using \eqref{eq:LG-Fexact} yields
\begin{equation}\label{eq:LG-Fzero}
\bra{\mathrm{phys}}\,F^a(x)\,\ket{\mathrm{phys}}=0
\qquad\Rightarrow\qquad
\bra{\mathrm{phys}}\,\partial^\mu W_\mu^a(x)\,\ket{\mathrm{phys}}
=
-\xi m_W\,\bra{\mathrm{phys}}\,\pi^a(x)\,\ket{\mathrm{phys}}.
\end{equation}
Equation \eqref{eq:LG-Fzero} is the configuration-space form of the Ward/Slavnov--Taylor identity needed.

\subsection{From \texorpdfstring{$\partial\!\cdot\!W^a$}{div W} insertions to amputated amplitudes}
Let $\mathcal{G}^{a}_{\mu}(x;\{y\})$ denote a connected Green function with one $W^a_\mu(x)$ insertion and
other insertions collectively denoted by $\{y\}$ (all taken to correspond to physical external legs in
an $S$-matrix element). Fourier transforming the $x$ dependence gives
\begin{equation}
\widetilde{\mathcal{G}}^{a}_{\mu}(p;\{y\})=\int d^4x\;e^{ip\cdot x}\,\mathcal{G}^{a}_{\mu}(x;\{y\}).
\end{equation}
Then the Fourier transform of $\partial^\mu W_\mu^a$ corresponds to multiplication by $i p^\mu$:
\begin{equation}
\widetilde{\mathcal{G}}^{a}_{\partial\cdot W}(p;\{y\}) \equiv
\int d^4x\;e^{ip\cdot x}\,\langle 0|T\{\partial^\mu W_\mu^a(x)\cdots\}|0\rangle
= i p^\mu \widetilde{\mathcal{G}}^{a}_{\mu}(p;\{y\}).
\end{equation}
Likewise define $\widetilde{\mathcal{G}}^{a}_{\pi}(p;\{y\})$ with a $\pi^a$ insertion.

Equation \eqref{eq:LG-Fzero} implies, at the level of such Green functions with only physical external
legs,
\begin{equation}\label{eq:LG-green}
i p^\mu \widetilde{\mathcal{G}}^{a}_{\mu}(p;\{y\}) = -\xi m_W\,\widetilde{\mathcal{G}}^{a}_{\pi}(p;\{y\}).
\end{equation}

Now apply LSZ reduction. For an external massive vector boson, the amputated object that becomes the
scattering amplitude is obtained by multiplying by $(p^2-m_W^2)$ and taking the on-shell limit; for an
external scalar Goldstone, one multiplies by $(p^2-\xi m_W^2)$ (because in $R_\xi$ gauge the Goldstone
propagator pole is at $p^2=\xi m_W^2$).
Denote the amputated, on-shell objects by $\mathcal{M}_\mu(W^a(p),\dots)$ and
$\mathcal{M}(\pi^a(p),\dots)$. The amputated version of \eqref{eq:LG-green} becomes the standard
Slavnov--Taylor relation
\begin{equation}\label{eq:LG-ST}
p^\mu\,\mathcal{M}_\mu(W^a(p),\dots) = i\,m_W\,\mathcal{M}(\pi^a(p),\dots),
\end{equation}
where we set $\xi=1$ in the final on-shell relation so that the Goldstone pole coincides with $m_W$
('t~Hooft--Feynman gauge). More general $\xi$ versions differ by $\xi$-dependent normalization factors
that cancel from physical $S$-matrix elements; the high-energy equivalence statement below is unchanged
in its standing-relevant content.

\subsection{Longitudinal polarization and the high-energy limit}
For an on-shell massive vector with momentum $p^\mu=(E,\mathbf{p})$ and $p^2=m_W^2$, a longitudinal
polarization can be chosen as
\begin{equation}\label{eq:LG-epsL}
\epsilon_L^\mu(p)=\frac{1}{m_W}\left(|\mathbf{p}|,\,E\,\hat{\mathbf{p}}\right),
\qquad \hat{\mathbf{p}}=\mathbf{p}/|\mathbf{p}|.
\end{equation}
One checks directly that $p\cdot \epsilon_L=0$ and $\epsilon_L^2=-1$ (Minkowski signature).
In the high-energy regime $E\gg m_W$, we have $E=|\mathbf{p}|+\cO(m_W^2/E)$, so
\begin{equation}\label{eq:LG-epsapprox}
\epsilon_L^\mu(p)=\frac{p^\mu}{m_W}+\cO\!\left(\frac{m_W}{E}\right).
\end{equation}

\subsection{Equivalence theorem (tree level, audit form)}
Combine \eqref{eq:LG-ST} with \eqref{eq:LG-epsapprox}:
\begin{align}
\mathcal{M}(W_L^a(p),\dots)
&= \epsilon_{L}^{\mu}(p)\,\mathcal{M}_\mu(W^a(p),\dots)\nonumber\\
&=\left(\frac{p^\mu}{m_W}+\cO\!\left(\frac{m_W}{E}\right)\right)\mathcal{M}_\mu(W^a(p),\dots)\nonumber\\
&=\frac{1}{m_W}\left(i m_W\,\mathcal{M}(\pi^a(p),\dots)\right)+\cO\!\left(\frac{m_W}{E}\right)\nonumber\\
&= i\,\mathcal{M}(\pi^a(p),\dots)+\cO\!\left(\frac{m_W}{E}\right).
\end{align}
The overall phase $i$ depends only on the convention used to define the scalar vs.\ vector LSZ
normalizations; at tree level and with consistent conventions it may be absorbed so that the
standing-relevant statement is:
\begin{equation}\label{eq:LG-final}
\mathcal{M}(W_L^a(p),\dots)
=
\mathcal{M}(\pi^a(p),\dots)
+\cO\!\left(\frac{m_W}{E}\right).
\end{equation}
Equation \eqref{eq:LG-final} is precisely the amplitude-class identification used in Lemma~\ref{lem:LG}.



\section*{References}
\begin{thebibliography}{99}

\bibitem{PeskinSchroeder}
M.~E.~Peskin and D.~V.~Schroeder,
\emph{An Introduction to Quantum Field Theory},
Addison--Wesley (1995).

\bibitem{WeinbergQTF1}
S.~Weinberg,
\emph{The Quantum Theory of Fields, Vol.\ 1: Foundations},
Cambridge University Press (1995).

\bibitem{WeinbergQTF2}
S.~Weinberg,
\emph{The Quantum Theory of Fields, Vol.\ 2: Modern Applications},
Cambridge University Press (1996).

\bibitem{Bertlmann1996}
R.~A.~Bertlmann,
\emph{Anomalies in Quantum Field Theory},
Oxford University Press (1996).

\bibitem{Witten1982}
E.~Witten,
``An $\SU(2)$ anomaly,''
\emph{Phys.\ Lett.\ B} \textbf{117} (1982) 324--328.

\bibitem{Fujikawa1979}
K.~Fujikawa,
``Path-Integral Measure for Gauge-Invariant Fermion Theories,''
\emph{Phys.\ Rev.\ Lett.} \textbf{42} (1979) 1195--1198.

\bibitem{Fujikawa1980}
K.~Fujikawa,
``Path Integral for Gauge Theories with Fermions,''
\emph{Phys.\ Rev.\ D} \textbf{21} (1980) 2848--2858.

\bibitem{Cornwall1974}
J.~M.~Cornwall, D.~N.~Levin and G.~Tiktopoulos,
``Derivation of Gauge Invariance from High-Energy Unitarity Bounds on the $S$ Matrix,''
\emph{Phys.\ Rev.\ D} \textbf{10} (1974) 1145--1167;
erratum \textbf{11} (1975) 972.

\bibitem{Chanowitz1985}
M.~S.~Chanowitz and M.~K.~Gaillard,
``The TeV Physics of Strongly Interacting $W$'s and $Z$'s,''
\emph{Nucl.\ Phys.\ B} \textbf{261} (1985) 379--431.

\end{thebibliography}


\end{document}

